\chapter{Entry}
\label{chap:entry}

\begin{chapteressay}
An event is not yet a fact. A spark in a wire, a bright point in a telescope, a
pea counted in a monastery garden, a wedge mark pressed into clay, a pulse in a
radio receiver: each may have happened, but happening is not the same as entering
the public record. Entry is the discipline by which an occurrence becomes a mark
someone else can inspect.

That distinction sounds small until one tries to write history from it. Most
histories of science can move quickly from an occurrence to a discovery because
the later record has already done the hard work. The table is printed. The plate
has been archived. The notebook has been edited. The instrument has acquired a
name and a lineage. In the moment itself, however, the occurrence is usually
more fragile. It flickers, is missed, is seen by only one person, appears once
and then not again, arrives without a stable unit, or is mixed with other marks
that could have been error. Before explanation, before law, before prediction,
there is the quieter discipline of asking whether anything has actually become
available for later use.

The first question, then, is not who was right. It is how an observation became
stable enough for rightness and wrongness to attach to it later. A private
sensation, a shouted report, and a remembered impression may all begin in
contact with something real, but none of them is yet a scientific record. The
mark has to pass through a rule for noticing, a rule for exclusion, a carrier
that can survive the moment, and some ordinary discipline by which another
worker can understand what would have counted and what would have been refused.

There is no insult to reality in saying this. The point is almost the opposite.
If the world is to correct a theory, the correction has to arrive somewhere. It
has to press on a scale, darken a plate, move a needle, interrupt a receiver,
fail to appear in a control, or leave a pattern in a table. Without such an
entry point, realism becomes theatrical: everyone gestures toward the world, but
no one has to say where the world touched the record. Entry makes the contact
less romantic and more exact. It asks for the channel, the threshold, the
material carrier, the competent observer, and the admissibility rule. Only then
can the later argument decide whether the mark was trustworthy.

The word ``ordinary'' matters. Entry is often made to look miraculous in
retrospect because later theory gives it a glow. The first wedge mark on clay,
the ruled table of planetary positions, the sketch beside a telescopic date,
the plate drying in a darkroom, the click of a counter, the tick of a marine
chronometer, and the crossed-out line in a laboratory notebook are not glamorous
objects by themselves. They are administrative, repetitive, and vulnerable to
mistake. But they are also the places where science becomes more than an
experience. A later reader can hold them still. A later critic can ask whether
the unit was stable, whether the category was legitimate, whether the absence
was recorded, whether the apparatus was competent, whether the mark was
transported without being silently changed.

The order of examples follows a problem rather than a chronology. Clay tablets
turn transactions into counts. Tycho's observatory and Mendel's garden turn
repeated acts into tables that outlive the people who made them. Galileo's
telescope and Hooke's microscope force sight to become dated drawings, printed
plates, labels, and controversy. Berkeley's criticism of Galilean abstraction
keeps a mathematical simplification from passing too easily as a physical
object. Marconi's wireless signal can be recognized only against static,
expectation, code, and silence. Chaitin's finite programs expose a limit inside
encoding itself. Photographic plates, Geiger counters, Faraday's notebooks,
Harrison's H4 chronometer, and Kelvin's failed age of the Earth all return to
the same entry problem: an occurrence is not yet the recorded form under which
the occurrence can be used.

The movement names are \lean{Fact}, \lean{DISTINGUISHABLE},
\lean{ADMISSIBLE}, \lean{COUNTABLE}, and \lean{ENCODED}. They are not a
chronology, and they are not offered as a hidden metaphysics behind history.
They name a minimum route. Something must first be treated as a candidate fact,
but candidacy is cheap. It must then be distinguishable from neighboring marks:
this pulse rather than that static, this trait rather than that ambiguous plant,
this star position rather than a smudge of atmosphere or instrument. It must be
admissible under a rule that can be stated without depending entirely on the
authority of the observer. It must be countable when the episode asks for
quantity, or at least capable of being placed in an order that makes comparison
public. Finally it must be encoded, carried, or stored in some medium that lets
it travel beyond the hands and eyes that first met it.

None of these steps is pure. A distinction can be trained into an observer and
then lost when the observer leaves. An admissibility rule can exclude a real
event because the apparatus was built for a narrower world. A count can make
comparison possible while hiding damaged categories. Encoding can preserve a
mark while stripping away the local conditions that made the mark intelligible.
Entry is therefore not a triumphal gate through which facts march cleanly into
science. It is a working threshold, full of tradeoffs. The public record gains
power because it can be inspected by people who were not present, but it also
becomes dangerous if the record forgets the conditions under which the mark
entered.

For that reason, the opening movement is also a theory of humility. It does not
promise a machine that turns experience into certainty. It describes the
discipline by which uncertainty becomes discussable. A good entry can still be
wrong, but it is wrong in public, with enough structure for correction to find
it. A bad entry may be attached to a true occurrence, but it leaves later
workers with nothing firm to inspect. The difference between those two cases is
not philosophical ornament. It is the difference between an event that can join
science and an event that remains stranded in testimony.

Failure is evidence here, not decoration. Kelvin's age of the Earth was not
merely a wrong number waiting to be corrected by later physics. It was a
disciplined entry that carried too little of the planet's thermal and
radioactive story. Blondlot's N-rays, in a later chapter, will show another
entry failure, where the mark dissolves when the apparatus of distinction is
tightened. The more basic lesson is that science must learn how to keep a
refusal. A null result, a missing signal, an ambiguous plate, an excluded
specimen, or a failed carrier may matter as much as a positive mark because it
tells later work where the apparatus stopped licensing claims.

The practical question underneath every section is therefore simple: what has
to be true of the record before someone absent can use it? Not what has to be
true of nature in some final sense, and not what a later textbook will make the
episode mean. The absent reader is the test. If the absent reader cannot tell
what was noticed, how it was distinguished, why it was admitted, how it was
counted or ordered, and how it was carried, then the event may have happened but
has not yet entered experimental science. If the absent reader can reconstruct
the relevant discipline, the mark has acquired a second life.

That second life is the beginning of experimental realism. An admitted mark is
not already true. The tougher claim is that the world can put pressure on
records only when records are built well enough to receive pressure. Bad entry
makes everything too easy: any impression can be called discovery, any silence
can be ignored, any unit can be counted, any diagram can masquerade as sight.
Good entry is harder and more modest. It allows the record to say yes, no, not
yet, not in this medium, not under this rule, and not with this apparatus. It
lets the experimenter be corrected by a mark that survives the experimenter.

Entry closes only when marks can travel. By then the record has moved through
ledgers, tables, images, notebooks, codes, instruments, and failed premises,
but the underlying question has not changed: how does an occurrence become
usable without pretending that the carrier is the world itself? Once that is
possible, a new problem appears. Marks accumulate, and what they cannot absorb
begins to persist beside them. Residue begins there.
\end{chapteressay}

\section{Number Refuses Possession}
\episodecoord{sixth century BCE / 1900 | Entry :: Fact -> DISTINGUISHABLE}

A monochord is a single string stretched over a measured wooden body and
held under steady tension. The full length of the string sounds first. That
open tone is the reference. A bridge --- a small wedge --- is then pushed
under the string at a chosen point along the body. The bridge sets a
new endpoint for the vibrating length: only the segment between the bridge
and one fixed end is the active sounding length compared against the open
string. The leftover segment is not the interval partner. The same string
is in play throughout, and the same tension. The only variable is the active
length, which the ruler reads directly off the body.

When the bridge is placed at half the body's length, the active section is
$L/2$. Plucked, it sounds an octave above the open reference. Its frequency
is twice the open string's, because pitch on a uniform string under fixed
tension is inversely proportional to active length. The length ratio,
active to open, is $1{:}2$; the frequency ratio is $2{:}1$.

When the bridge sits two-thirds of the way along the body, the active
section is $2L/3$. Plucked, it sounds a fifth above the open reference; the
frequency ratio is $3{:}2$. At three-quarters, the active section is
$3L/4$, a fourth above the open reference; the frequency ratio is $4{:}3$.

The ear hears the frequency ratios as consonance: the octave is
unmistakable, and the fifth and fourth are stable consonances in the same
practice. The ruler reads the inverse length ratios directly off the body.
The same small whole numbers survive both checks. Another person can set the bridge
at the same marks, pluck, and hear the same intervals. The check owes
nothing to Pythagorean doctrine, to the player's skill, or to any private
claim about cosmic harmony. It is a public observation made in two media at once.

Then the diagonal of the unit square breaks the confidence. The diagonal
is constructed as plainly as the side: a square is drawn, the
corner-to-corner line is laid down with a straightedge. The construction
yields a definite length, and that length is the square root of two. But no
ratio of whole numbers $p/q$ equals that length. Integer ratio stops being a
universal medium for entry. Mathematics has not failed; the proof that the
square root of two is irrational is one of the cleanest results in
classical number theory. What has reached a limit is the carrying medium.
The geometric record can construct a magnitude that the integer-ratio entry
system cannot inscribe. The record meets that magnitude and admits the
refusal rather than pretending to carry what it cannot.

\begin{phenom}{The Pythagoras--Planck Entry Effect~\cite{euclid300bc,planck1900}}
\label{ph:pythagoras-planck-entry}

\PhStatement
A mark enters science only when the medium that carries it can also name what
it cannot carry.

\PhOrigin
Euclid's account of incommensurable magnitude preserves the old Greek refusal
of integer ratio as universal carrier. Planck's blackbody work, his
study of radiation from an ideal heated body, later forced physicists to
accept that energy exchange between matter and light could only proceed in
finite packets, not in arbitrarily small amounts. The historical
settings are different; both moments discipline entry by refusing unlimited
continuous refinement.

\PhObservation
Ratio, construction, and measurement need not agree. A construction can make a
magnitude public while the available numerical carrier fails to encode it
without residue.

\PhConstraint
The record distinguishes the constructed mark from the carrier by which the
mark travels. Carrier failure is not permission to erase the mark or to
pretend the carrier was sufficient.

\PhConsequence
Entry begins with a double obligation: admit the mark, and preserve the refusal
that tells later instruments where the first carrier stopped.
\end{phenom}

\section{Counting as Public Order}
\episodecoord{ancient tokens / 1889 | Entry :: COUNTABLE -> ENCODED}

Kushim is the name --- or perhaps the title of an office --- read on an early
Uruk administrative tablet from around 3000 BCE. The tablet records a
transaction in barley, and the sign-group read as Kushim survives because
that transaction needed a written tally. The record has preserved an entry,
not a biography. Whether Kushim names a person, an office, or a scribal
responsibility, the public fact is not a life. It is a ledger mark.

Counting is often treated as the least mysterious part of science, as if
numbers arrive after the objects have already been cleanly given. The
earliest ledgers teach the reverse. A count imposes a discipline on what may
be entered as a unit. A sheep, a jar, a loaf, a measure of barley, a day of
labor, a pulse, a pea, a click, or a collision is countable only after some
practice has decided what separates one from another.

The clay tablet has every feature of an entry apparatus, even though its
work is too familiar to look like instrument-making. It receives a mark. It
stabilizes the mark after the transaction is over. It can be checked by
someone who was not present. It carries omissions and exclusions as surely
as it carries units. It can be broken, misread, forged, corrected, filed,
or lost. The scribe needs no theory of natural number to make this apparatus
work; the scribe needs a public enough rule: this kind of thing, under this
administrative category, at this point in the transaction, is one item for
this ledger.

The phrase ``for this ledger'' keeps the episode realistic. The ancient
count is not a universal claim about what barley is; it is a local discipline. Barley counted
for storage is not the same as barley tasted, planted, cooked, spoiled,
taxed, or measured for ritual distribution. The grain is not false in any
of those cases, but the unit is different. A unit is never simply found.
It is prepared by a practice that decides what differences matter and what
differences will be ignored. Moisture, quality, labor, ownership, and
spoilage may matter elsewhere and disappear here. The tablet is powerful
because it makes a transaction survive; it is limited because survival
requires reduction.

No line of influence runs from an Uruk account tablet to Peano's axioms.
Giuseppe Peano, in 1889, isolated the discipline that counting practice
already used and presented it in purified form: a starting point, a
successor operation, identity conditions that prevent collapse, and
induction, the rule by which a property holding at one step is carried
through to the next. The clay tablet leaves those conditions implicit. It
uses a rough administrative version of them.

Each new mark advances the state of the record. The ledger registers not
only that an item was recognized, but that the recognition has been carried
forward. The public count is therefore temporal as well as numerical. It has
a before and an after. It has a rule by which the next entry stands as next
rather than as a correction, a duplicate, or an unrelated mark. If that rule
is unstable, counting collapses into inscription without order.

Peano's formalism strips the goods away. The successor of a number is again
a number; two distinct numbers cannot share the same successor; the first
element is itself no one's successor; and induction lets the sequence carry
a property through all later entries. Notation and convention vary, but the
discipline is clear enough: counting is an ordered public procedure, not a
heap of names. It tells the record when one more has actually been added.

Experimental science relearns this lesson. Mendel cannot count a wrinkled
pea until ``wrinkled'' is a stable category. A cloud chamber track cannot
be counted until it can be distinguished from background. A Geiger counter
click cannot be counted until the apparatus has dead time, threshold, and
calibration. A clinical trial cannot count a case until the protocol has
decided what qualifies. In each
setting, arithmetic is the visible surface of a deeper entry discipline.
The equation may contain a number, but the experiment had first to decide
what the number was allowed to count.

Counting is not the same as seeing many things. A pile is not a count. A
crowd is not a count. A noisy trace is not a count. A count begins when the
record has a unit, an increment rule, and a refusal rule. The refusal rule
is as important as the increment rule. A damaged jar may not count. A
duplicate receipt may not count. A pea whose trait cannot be assigned may
not count. A pulse below threshold may not count. Those refusals are not
embarrassments outside the evidence. They are the conditions under which
the evidence is allowed to become numerical.

A table of counts looks public because anyone can inspect the totals. The
totals are only as public as the entry procedure that produced them. If a
later reader cannot tell which marks were admitted, what counted as one,
whether duplicates were removed, how ambiguous cases were handled, and
where the ledger began and ended, the number is not public --- it is a
number whose entry rule has been hidden from view. Hidden practice is not
always corrupt; often it is simply skilled or ordinary. The danger
is that science needs the count to survive the person who made it.

The Kushim tablet survives because administration had to outlast presence.
Delivering parties left. Receiving parties left. Officials checking the
account could arrive later. The tablet made a transaction portable across
that gap, though never fully transparent. It made enough of the transaction
inspectable for a later administrative act. Experimental entry inherits the
same minimal ambition: not total access to an event, but a disciplined
object strong enough to support later use.

The tablet is sometimes described as turning reality into data. The tablet
turns a selected transaction into an entered mark. Data, in the strong
experimental sense, comes later, only when many such marks can be organized,
compared, refused, and transported through additional instruments. Before
the world can become a dataset, something has to be admitted as a countable
unit. The unit is the first compromise between occurrence and record.

The formal side carries its own warning. Once counting is formalized, the
formalism carries whatever was admitted under the unit rule. The successor
operation will preserve order regardless of
whether the scribe counted barley measures fairly, whether Mendel's
categories were stable, whether a detector double-counted, or whether a
dataset silently excluded failed trials. Formal arithmetic preserves order;
it cannot, by itself, certify the entry conditions. That certification
belongs to the apparatus, the protocol, and the source check.

A theory may compute from a number, but the number is not real because it
appears in a calculation. It becomes available as a claim
about reality only when its entry can be inspected --- when the number is tied
back to a unit, a threshold, a category, a carrier, and a refusal rule.
Otherwise the computation may be exact while the count remains unlicensed.
The danger is not mathematical error but perfect arithmetic attached to an
unexamined ledger.

The stop rule is part of the same discipline. A count without a boundary is
not yet a count; it is an activity. The ancient account has to know where
the transaction ends. The experimental run has to know when the exposure
is over, when the sample is closed, when a pulse belongs to this interval
rather than the next, and when a correction is still a correction rather
than a new entry. The stop rule is easy to hide because it often looks
practical rather than conceptual. Yet it decides the denominator, the rate,
the table row, and the possibility of comparison. Without it, ``one more''
has no stable place to land.

A private count can be remembered. A public count needs to be checkable.
The clay mark, the column total, the tally sheet, the photographic
catalogue, and the detector log all make recounting possible. Recounting is
not mere suspicion that the first worker may have erred --- it is a
condition of public use. A count that cannot be recounted, audited, or at
least traced back through its entry decisions asks later workers to inherit
trust without the means of criticism. The ledger form turns trust into a
practice: the entries
are visible, the path to the total is visible, and a later reader can begin
to object.

The difficult cases tell why this is not bookkeeping trivia. One unit may
split into two. Two marks may turn out to be the same object entered twice.
A specimen may change category after closer inspection. A pulse may arrive
during detector recovery, or an image may show a mark at the edge of
visibility. The successor operation says how to advance once a unit is
admitted. It cannot resolve those questions. They belong to the apparatus
and the protocol. They are the places where counting touches matter, habit,
fatigue, craft, and institutional pressure.

A realistic history also notices that counts become social facts before
they become scientific facts. The administrative tablet was useful because
people could act on it: store the goods, settle the account, assign
responsibility, or remember an obligation. Modern experiments retain this
social edge. A published count can allocate credit, end a dispute, justify
an instrument, close a trial, or open a new field. That social force is not
what makes a count false; it makes the entry conditions more important.
The more action a count can support, the more carefully the record needs
to show how the count was made.

Category drift follows from this. A ledger accumulates units only if the
same entry rule remains stable enough across time. But practices change ---
assistants learn, instruments are repaired, protocols are clarified, samples
age, categories are renamed, old marks are read under new interests. The
count then carries two histories: the history of the things counted and the
history of the rule by which they were counted. A realistic account keeps
both visible. If the rule drifts silently, the total may look continuous
while the units have changed underneath. If the rule is annotated, the
drift becomes part of the evidence.

Every count reduces. It leaves out texture, context, and difference in
order to preserve order. That reduction is necessary for comparison. The mistake is pretending that
no reduction occurred. A tablet
that counts barley measures is not a poem about barley. A detector count
is not a biography of every particle interaction. A table of offspring
traits is not the lived development of each plant.
A count is a compressed public object. It earns its evidential force by
exposing enough of the compression for later workers to know what was lost.

The absent reader, again, is the test. Imagine someone who trusts neither
the scribe nor the experimenter personally. That reader can still use the
record if the unit rule, the successor rule, the stop rule, and the
refusal rule remain recoverable. The reader needs no omniscience. The
reader needs a disciplined path from mark to count. Objectivity here is
not detachment from any particular standpoint; it is a record arranged so
that someone elsewhere can perform a limited, critical reconstruction.

Both records are incomplete alone. The tablet without formal discipline
preserves an entry while leaving its logic implicit. The formal system
without entry discipline describes counting while saying nothing about what
was counted. Experimental science needs both: the ledger ordered enough to
travel, and the order answerable to the material practice that filled it.

The transition from \lean{COUNTABLE} to \lean{ENCODED} matters because a
count that cannot be carried is only a local performance. Encoding gives
the count a durable body: a notch, a token, a wedge, a column, a printed
numeral, a digital register, a database row. Encoding also hardens the
decision. Once the mark has been stored, later readers may treat it as if
the unit had always been obvious. The encoded count begins to look like
the event itself. The distinction has to be kept simple: the event
happened, perhaps; the unit was admitted, under a rule; the count
advanced; the encoded mark survived.

These rules are permanent disciplines of science, not preparatory to it.
The most advanced detector still decides what increments the tally; the
most elegant model inherits the entry rule of its data; the most persuasive
plot depends on whether each point deserved to be there. Peano and Kushim
together say that a counted fact is not just a fact plus a number: it is
an occurrence that has passed through a public order of units, successors,
and refusals.

The same discipline reappears in later materials. Tycho's tables need
stable rows. Mendel's ratios need stable traits. Galileo and Hooke need
visual marks that count as appearances rather than marvels. A photographic
plate needs conventions for identifying a trace. A Geiger counter needs a
rule for when a click counts as one event. Faraday's notebooks need dated
entries that keep trials from collapsing into anecdote. Kelvin's age
calculation shows what happens when a record is disciplined in its
arithmetic but incomplete in its entry premises. The ledger form remains
inside later experimental records.

\lean{COUNTABLE} is not a claim that the world is naturally parceled into
neat units. It is a claim about a record that has earned the right to
advance by one. That right is practical, fragile, and revisable. It can be
challenged by damaged material, better instruments, new categories, or a
later audit. Once earned, the count is a small public machine for carrying
sameness through time. Kushim gives that machine its clay body. Peano gives
it its exposed logic. Experimental science needs both --- because it keeps
asking not only how many, but how the record was allowed to say one more.
Without that answer, the count is only a number, with no claim on the world
it appears to describe.

\begin{phenom}{The Peano--Kushim Ledger Effect~\cite{peano1889,schmandtbesserat1992}}
\label{ph:peano-kushim-ledger}

\PhStatement
Counting makes an occurrence public by admitting it as a unit that can advance
a durable ledger.

\PhOrigin
The early administrative tablet associated with Kushim preserves a transaction
through marks that can be inspected after the event. Peano's axioms later
isolate the successor discipline that lets one entry follow another without
collapse. The two settings are not historically continuous in the simple
sense, but they expose the same constraint: a public count needs units, order,
and a rule for advancement.

\PhObservation
The tablet preserves selected units under a rule, not the whole harvest, the
labor, the weather, or the private intention of the scribe. Peano's formal
order selects no units; it states the conditions in force once something has
been admitted as countable.

\PhConstraint
An experimental count identifies what increments the ledger, what remains
outside it, and how ambiguous cases are handled. A unit that cannot be
distinguished under the counting rule is not yet a unit for this record. A
formal successor can carry an unlicensed unit, but it cannot make the unit
licensed.

\PhConsequence
Counting is not arithmetic added to science after the fact. It is the first
public order by which later comparison becomes possible. Every later table,
detector, trial, and database inherits this condition: before a number can
discipline theory, the record shows what was allowed to become one.
\end{phenom}

\section{Tables That Survive the Observer}
\episodecoord{late 1500s / 1866 | Entry :: COUNTABLE}

Tycho Brahe was a Danish astronomer working in the late sixteenth century,
before the telescope. From his observatory on the island of Hven and later
in Prague, he and his assistants used large angular instruments --- quadrants,
sextants, an armillary sphere --- to take repeated, calibrated position
measurements of stars and planets and to record each into dated tables.
Gregor Mendel was a nineteenth-century monk in Brünn who grew pea plants in
his monastery garden, performed controlled crosses between them, and
counted the traits of their offspring across multiple generations in
tabular form. Each turned a kind of seeing into a record that outlived the
seeing.

Tycho's observatory turned position into discipline. Mendel's garden turned
inheritance into discipline that lasted past the plants. Neither case is
merely ``a lot of data.'' The table is the apparatus after the apparatus.

A quadrant, sextant, armillary sphere, garden bed, pollen brush, or seed
packet can make an observation possible, but the table is what lets the
observation stay at work after the immediate apparatus has stopped acting.
Tycho's instruments were rebuilt, adjusted, criticized, and eventually
lost. Mendel's plants died, reproduced, were selected, and
disappeared into later seasons. What remained usable was not the original
sight of a star or the original feel of a pea, but an ordered record in
which a later worker could ask what had been entered, what had been
excluded, and how one entry stood beside another.

The table is easy to underestimate because it looks passive. In practice
each row is an active decision: this mark belongs in this row rather than
that one, under this heading rather than another, at this date, with this
correction. The blank cell is also meaningful. A missing observation, an
empty cell, a suspicious regularity, or a category with no entries can all
become part of the later argument. The table is not just storage. It is a
public shape for attention.

Before the telescope, astronomical precision depended on large instruments,
trained observers, repetition, geometry, and correction. A single sighting
could not carry the burden. Celestial position had to be turned into a
practice: align, read, record, repeat, compare, correct, and enter. The
observatory was not only a place where the sky was seen. It was a place
where seeing was routed into durable coordinates. The table made an
observation legible as part of a series.

Astronomical observation arrives without a ready-made table. Time scale,
angular calibration,
instrument alignment, assistant coordination, atmospheric conditions, and
the reconciliation of repeated readings each enter the record by decision.
Each decision can change the entry. Tycho's achievement was not that he
escaped mediation but that he built mediation strong enough for later
readers to find patterns Tycho himself missed. The table preserved more
than Tycho's cosmology.

A record can outlive the theory that first gathered it. Tycho's own
cosmological commitments left his observations available to later
astronomical work. The table carried positions in a form that could be
detached from some of the observer's interpretation, not because tables are
neutral but because a well-made table separates enough of the entry rule
from the explanatory claim that later workers can reuse, challenge, or
transform the record.
The recorded position survives its first explanation.

Mendel's garden teaches the same lesson in a different register. The
experiment begins not when a ratio appears but when a trait can be entered
in a stable enough way to be counted across generations. Smooth and
wrinkled, green and yellow, tall and short: such categories look simple
only after the experimental practice has performed its sorting. Plants are
selected. Crosses are controlled. Self-fertilization and hybridization are
managed. Offspring are assigned to categories. Counts are carried through
generations. The ratio comes late; the entry discipline comes first.

Mendel's paper is powerful because it gives the reader not merely
conclusions but a tabular order of inheritance. The plants are gone, but
their counted traits remain arranged for inspection. A later reader can ask
why the categories were chosen, whether ambiguous plants were excluded, how
the generations were separated, how large the counts were, and why some
ratios look so persuasive. Those questions are possible because the record
is not a single claim that heredity behaves this way; it is a structured
table of entries from which the claim is drawn.

That structure also makes the record vulnerable. Tables invite trust, but
they also invite audit. A table that looks too clean may become suspicious.
A table with hidden exclusions may mislead. A table with drifting
categories may combine unlike units. A table with unreported failures may
make a pattern look stronger than the experimental practice warrants. The
answer is to ask what discipline allowed each row to enter. The table is the beginning
of public criticism, not the end of it.

Both cases convert repetition into survivable order. Tycho's repeated
positions and Mendel's repeated crosses are not casual repetitions; they
are repetitions under a rule. The later
entry is supposed to be comparable to the earlier entry. If the instrument
has changed, the category has shifted, the assistant has used a different
rule, or the date has been miscopied, the table may still look orderly
while its comparability has broken. Repetition becomes scientific only when
the table makes the sameness of the entry procedure inspectable.

The connection to \lean{COUNTABLE} is direct. A table is a count extended
into two dimensions: not just one more, but one more under a heading, in a
row, at a time, in a series, beside other entries. Headings tell the reader
what sort of thing the count claims to count. Rows anchor entries in
ordered sequence. Totals compress entries while inviting a recount.
Footnotes, corrections, and blank cells mark where the simple count needed
help.

Tables can therefore be more realistic than anecdote. An anecdote may be
vivid, but it asks the reader to stay close to the witness. A table creates
distance. It lets the reader compare marks without needing to relive the
observer's whole experience. That distance is what gives an absent reader
room to ask whether the marks deserve the pattern being claimed for them.
Science needs that room because the observer is often sincere, skilled, and
still wrong.

The table also changes what counts as a mistake. In a single observation,
an error may vanish with the moment. In a table, error can propagate. A
copied number, a shifted column, a wrong date, a category entered under the
wrong heading, or a correction applied twice can travel far beyond the
original act. The same durability that carries error also allows error to
be found. A later reader can notice a row that breaks the sequence, a total
that fails to add, an outlier that demands explanation, or a blank where an
entry was expected. The table preserves both evidence and the possibility
of repair.

Tables cannot guarantee truth; they make truth and error more publicly
vulnerable. They give later work a record to return to. Kepler could use
astronomical positions against inherited circular expectations because the
positions had been preserved in a disciplined form. Later geneticists could
return to Mendel's counts, admire them, doubt them, reinterpret them, and
embed them in a different biological framework because the counts had
survived as an ordered record. The table lets a result have a future.

Cross-time readability is not automatic. Headings need to remain
interpretable. Units need to be known.
The provenance of entries needs to be strong enough to keep later readers
from guessing at the wrong practice. Source discipline belongs inside the
experiment rather than after it. A table without source discipline becomes
an image of order. It may still persuade, but persuasion is not entry.
Entry requires that the reader can reconstruct enough of the table's making
to know what its order licenses.

The table also makes selection visible. No table can contain everything.
Tycho had to choose which observations mattered, how to correct them, and
how to present them. Mendel had to choose traits that could remain discrete
enough for counting. Experiment narrows the world to make a public question
answerable.
The defect comes when narrowing is hidden. A table is honest when it lets
the reader see both its power and its frame.

The frame includes the labor behind the table. Tycho's observatory depended
on large angular instruments, masonry, metalwork, assistants, night
routines, and corrections, not only on a brilliant eye. Mendel's garden
depended on selected pea varieties, controlled pollination, bagging and
handling of flowers, seasonal patience, and recordkeeping, not only on a
sudden insight about inheritance. The finished table can make that labor
disappear because the entries look abstract. A row is a condensed history
of work. A column heading is a stabilized practice. A total is a social
agreement that these marks may be gathered.

The table is objective in practice: built from procedures that can be
inspected, repeated, and criticized. The table is
not outside human work; it is human work disciplined so that the mark can
survive the worker. The observer can leave and the record can continue
because the mediation has been made durable enough to be tested by others.

Tycho--Mendel also clarifies the difference between accumulation and entry.
Observations can be piled up without producing a scientific table. A
notebook full of impressions is not yet an ordered record. A drawer full of
plates is not yet a catalogue. A garden full of plants is not yet heredity
data. The entries need to be placed so that relations among them become
inspectable. The table gives pattern a public place to appear.

Once the pattern appears, it can exceed its maker. Tycho's positions
mattered without him being right about the architecture of the heavens.
Mendel's counts mattered before chromosome theory could explain them.
The table allowed their records to become partially independent of their
local interpretations. That independence is never total, but it is enough
to let history move. A mark that survives as part of a table can enter
arguments its maker never foresaw.

Correction is one reason this independence is possible. A raw sighting of
a planet has almost no future unless the reader knows how to place it: time,
reference point, instrument, angular reading, and correction practice. A
raw plant description has almost no future unless the reader knows which
generation, cross, trait category, and exclusion rule produced it. The table lets these corrections become
part of the record's grammar. It says not merely what was seen but how the
seen mark was made fit to stand beside other marks.

That fitting is never innocent. The table smooths the world enough to make
comparison possible. It turns an evening of observation into a row and a
season of cultivation into a count. It forgets weather, fatigue, failed
attempts, hesitations, small repairs, and much of the bodily labor that
made entry possible. This forgetting can be legitimate only if the table
keeps enough traces of its own terms. The reader needs to know what kind of
forgetting the table has performed, not every detail of the night sky or
the garden bed.

The two examples sharpen each other. Astronomical position looks
quantitative from the start, yet it rests on embodied judgment: aligning
sights, reading scales, deciding when an observation is good enough, and
correcting for known disturbances. Mendelian inheritance looks categorical
from the start, yet it becomes numerical only after a biological continuum
has been sorted into repeatable classes. The first case reminds us that
numbers come from instruments. The second reminds us that categories come
from practices. The table needs both.

The strongest table is not the neatest. A suspicious table may have rows
that are too perfectly regular, totals that hide exclusions, or categories
that accepted only friendly cases. A trustworthy table may include
awkwardness: failed entries, corrections, uneven sample sizes,
qualifications, or notes that keep a future reader from mistaking the form
for the event. Realism often enters through these unpretty marks. They show
that the table was made in contact with resistant material rather than
written backward from a desired conclusion.

The blank cell deserves attention because it can mean nothing was seen,
nothing was measured, nothing was entered, the apparatus was unavailable,
the category was inapplicable, or the compiler chose not to report. Those
are different absences. A mature table teaches the reader which absence is
in play. In astronomy, an omitted observation may mark weather, scheduling,
instrument failure, or irrelevance. In breeding experiments, an absent
category may mark a biological result, an exclusion, or a failed cross.
The blank is not empty until the entry rule says what kind of emptiness it
is.

The same is true of totals. A total feels like closure, but it is often a
compressed argument. It says that these entries belong together and that
their differences may be temporarily set aside for this purpose. Tycho's
positions cannot simply be added, but they gather into series that
constrain orbits. Mendel's offspring can be added, but only after the
offspring have been sorted into categories. In both cases, the table
creates a public many. The total or series is not a pile; it is a claim
that the entries share enough of a rule to be read together.

From entry the table turns toward residue. Once observations are tabled,
their leftovers become harder to ignore. An outlier remains in its row. A
missing case leaves a gap. A correction leaves a trace. A ratio that is too
clean or too rough becomes a question rather than a rumor. The table gives
residue a place to accumulate. It lets future work see not only what the
record claimed, but where the claim was strained.

Tables can be misused. The same format that exposes a bad total can make a
weak pattern look authoritative. The same rows that invite audit can
intimidate readers who mistake tabular order for settled truth. The table
is a public instrument, and public instruments can be used well or badly.
Their realism lies in the possibility of inspection, not in the appearance
of order.

The table extends the ledger effect rather than replacing it. Peano--Kushim
showed how one more becomes public. Tycho--Mendel shows how many ones
become a structure. The table is a disciplined many: a record in which each
entry keeps enough identity to be counted and enough relation to be
compared. Entry now has more than a mark. It has order.

That order is why later science can become cumulative without becoming
merely repetitive. The next worker can begin from an ordered residue of
earlier work, then test, correct, extend, or refuse it. This is not
secondhand science in the weak sense; it is the normal public life of
science: firsthand marks arranged so that responsible secondhand use
becomes possible. The table turns observation into something that can have
heirs. Those heirs may be loyal, skeptical, or transformative. The table
gives them a record determinate enough to inherit.

\begin{phenom}{The Tycho--Mendel Table Effect~\cite{mendel1866,tycho1602}}
\label{ph:tycho-mendel-table}

\PhStatement
A table is an instrument for making repeated observation survive the observer
as ordered, comparable entry.

\PhOrigin
Tycho's astronomical program depended on sustained tabulation of celestial
positions from large naked-eye instruments and repeated readings. Mendel's pea
experiments depended on categorical counts across generations, controlled
crosses, and traits stable enough to enter columns. The episodes are
historically unrelated in the simple genealogical sense, but both produced
records whose force came from ordered tables rather than from a single act of
seeing.

\PhObservation
The table carries repeated marks in a form that can be sorted, checked, and
re-read. It also exposes missing rows, ambiguous categories, unstable headings,
wrong totals, and suspiciously neat patterns. The table preserves a series
while giving later readers a surface for criticism.

\PhConstraint
Tabular evidence is admissible only when the rule for entering rows and columns
is public enough for later readers to understand what each count could have
excluded. A table that hides its units, corrections, stop rules, or category
choices may look ordered without making its order experimental.

\PhConsequence
The observer leaves, but the table remains. Experimental science becomes
cumulative only after marks acquire this second life: an order strong enough to
outlive the observer and vulnerable enough to be challenged by someone else.
\end{phenom}

\section{Seeing Enters a Record}
\episodecoord{1610 / 1665 | Entry :: Fact -> ADMISSIBLE -> ENCODED}

In 1610 Galileo Galilei published \emph{Sidereus Nuncius}, recording what
his telescope had shown him: small bodies moving night by night near
Jupiter, rough mountainous shadows on the Moon, and stars in places where
ordinary sight had seen only a pale band. In 1665 Robert Hooke published
\emph{Micrographia}, recording what his microscope had shown him: the
compartments of cork, the armored surface of a flea, and the rough point
of a needle. Each book gave the new
instrumented sight a public form --- dated drawings, copperplate
engravings, descriptions, and explicit scale conventions --- through which readers far
from the instrument could argue about what was seen.

The telescope made no one right by magnifying the sky. It made new
argument possible by forcing sight into notebooks, drawings, dates, names,
and public controversy. Hooke's microscope had the same effect in another
direction. Magnification without inscription would have remained marvel.
Magnification with inscription became an entry problem.

The shift these episodes mark is not that the human eye suddenly became
more powerful. A stronger eye changes nothing public if its report cannot
travel. The shift is that instrumented sight had to learn how to leave a
mark. Telescopic sights entered public argument only when they were dated,
drawn, compared, narrated, and sent to readers who were not at the
eyepiece. The microscopic view entered science through plates,
descriptions, scale conventions, and a printed performance of care.

Image-making functions here as apparatus, not decoration. The drawing is
not an afterthought added to a finished observation. It is one of the
means by which the observation becomes admissible. A sketch beside a date
can preserve a changing relation among points of light. A printed
engraving can give distant readers a disciplined substitute for the
fleeting microscopic view. A label can tell the reader which part of the
image matters. A sequence can distinguish motion from accident. These are
not neutral acts; they are entry decisions.

Galileo's case is especially severe because the telescope itself had to
be made credible. A critic could always ask whether the instrument
produced illusions, whether the observer was trained, whether the points
were stars, whether the sequence was correctly dated, or whether the
report had outrun the evidence. The answer could not be simply ``I saw
it.'' Instrumented sight needs more than testimony because the instrument
changes the conditions of seeing. The record has to show how the new
sight can be inspected by others who may distrust both the instrument and
the observer.

The dated sequence of Jupiter's companions is a good example of entry
pressure. A single view could be dismissed as confusion, artifact, or
misplacement. A series of entries lets relation become visible. Points
appear on different sides, change positions, disappear, and return. The
record holds no modern orbital theory in hand yet, but it has made the
candidate fact harder to treat as a private impression. The sight has
been put into time and can be checked against a pattern.

In \emph{Sidereus Nuncius}, the small star symbols beside Jupiter carry
much of this work. They are spare, almost dry marks: a date, a line of
position, a few points. Their plainness is useful. The reader is not
asked to admire a marvel; the reader is asked to follow displacement
from night to night. The claim enters through a dated sequence before it
enters through later celestial mechanics.

The Moon's surface raises a different problem. Roughness and shadow have
to be translated into drawing. The observer has to decide how to render
contrast, which features to emphasize, and how to persuade a reader that
the mark belongs to topography rather than lens, eye, or imagination.
Here the drawing carries both evidence and interpretation. It is a
mediated record of sight. That is not a mark against it; it makes the
mediation visible enough to argue about. The public record gains force
because the drawing can be compared, criticized, copied, and repeated,
even though it is not the Moon.

Hooke intensifies the same problem because the microscope makes the
ordinary world strange. A flea becomes architectural. Cork becomes
cellular. A needle point becomes rough. The plate asks the reader to
accept not only that Hooke saw something, but that the instrumented view
has been translated into a proportionate, meaningful, and shareable
image. The printed plate is therefore both evidence and pedagogy. It
trains the reader in how to look at a world that cannot be inspected by
naked sight.

The copperplate engravings in \emph{Micrographia} are not casual
illustrations. They are the publication technology by which a fleeting,
trained view is made available to people without Hooke's microscope,
specimen, lamp, or hand. Cork's little compartments, the flea's armored
surface, and the rough needle point become public through a chain of
preparation, viewing, drawing, engraving, and printing. The chain is long
enough to be dangerous and concrete enough to be criticized.

The realism of these images lies in their discipline, not in any supposed
purity. Neither Galileo nor Hooke gives raw sight. They give selected,
prepared, and rhetorically arranged records of sight. The lens has
aberrations. The observer has habits. The drawing hand has conventions.
The engraver or printer introduces another transformation. The page
frames what the eye is meant to notice. All this mediation can be a
source of error, but it is also the only route by which the sight becomes
public. Refusing mediation would not give a cleaner fact; it would return
the event to private seeing.

\lean{ADMISSIBLE} is the central term here. An instrumented visual claim
is not admissible merely because it is vivid. Vividness can make error
more persuasive. The admissibility rule asks whether the record tells a
later reader enough about the conditions of seeing: what instrument was
used, what was selected, what was drawn, what was omitted, how the image
relates to a series, and what alternative explanations remain possible.
The rule is not absolute, but without some version of it magnification
becomes theater.

The record also has to protect against the authority of novelty. New
instruments often produce astonishment before they produce discipline. A
telescope or microscope can make observers feel they have gained direct
access to a hidden world. The danger is that astonishment may be mistaken
for entry. The first public question is less grand: what mark has the
instrument licensed, and how can that mark be carried? A marvelous sight
that cannot be dated, drawn, repeated, or otherwise translated remains
weak as evidence. A modest mark with a stable carrier can perform more
scientific work than a marvel.

Telescope and microscope together show that public seeing is social
before it becomes general. Readers have to learn the instrument's
authority. Makers have to build devices usable by more than one person.
Demonstrations have to separate skill from fraud and instrument from
illusion. Publications have to give enough instruction that others can
try, fail, object, and refine. The image alone is not enough. The image
belongs to a practice of witness-making; it recruits absent readers into
a community of trained observers.

A sequence of drawings is already a table in visual form. Galileo's
dated positions near Jupiter ask the reader to read images relationally,
not as isolated marvels. Hooke's plates gather specimens into a visual
catalogue whose order teaches comparison. The visual record adds a new
carrier to entry; entry can now move through line, shade, scale, and
printed plate.

Scale is especially important. A microscopic object cannot enter the
page at its ordinary size. It has to be enlarged, framed, and related to
a convention of measurement or comparison. The enlargement is useful
precisely because it is not the thing itself. It is a translation that
gives the reader access to relations otherwise unavailable. But the
translation has to declare enough of its terms. A huge flea on a page is
not evidence unless the reader knows that the size is an instrumented
representation. A magnified world without scale can become fantasy.

The same warning applies to telescopic images. A small dot in a sketch
may stand for a luminous body, but only within a rule of observation and
notation. The dot is not the body; it is the notation through which the
body enters the record. The image is not nature itself, but neither is
it merely invention. It is a disciplined carrier whose claim depends on
the instrument, the observer, the inscription, and the possibility of
later criticism.

Hooke's engravings make another difficulty plain. The more beautiful the
image, the easier it is for readers to forget the chain of translation.
Beauty can stabilize attention; it can also hide the labor that made the
mark. A plate may look inevitable once it is printed. In fact it has
passed through specimen preparation, lighting, lens work, observation,
drawing, engraving, printing, and readerly interpretation. Each step can
sharpen the image or distort it. Realism requires naming the chain, not
pretending the image simply fell out of the instrument.

Galileo's public controversy makes the same point by another path. If
everyone could immediately see and agree, inscription would matter less.
But agreement was slow to arrive. Some readers lacked instruments. Some
distrusted the device. Some could not see what Galileo reported. Some
interpreted the sights differently. The printed record became a way of
carrying the dispute without requiring every reader to be present at the
original eyepiece. The record made controversy portable rather than
ending it.

This portability carries sight from \lean{Fact} through \lean{ADMISSIBLE}
to \lean{ENCODED}. The sight first appears as a candidate fact: there is
something there. It becomes admissible when the record supplies enough
conditions for judgment: date, instrument, sequence, drawing,
description, and exclusion of obvious alternatives. It becomes encoded
when the sight travels as a mark in a notebook, printed page, plate, or
copy. At each step something is gained and something is lost. The
private intensity of seeing is lost; public inspectability is gained.

That loss is the trade. Private intensity is often exactly what science
has to surrender. The observer's astonishment may motivate the work, but
it cannot be the evidence by itself. The evidence becomes poorer in one
sense in order to become richer in another. It loses the fullness of the
moment and gains a carrier. It loses immediacy and gains a future. It
loses the observer's authority and gains the possibility of criticism by
people who were not there.

That surrender is difficult because seeing feels sovereign. A person who
has looked through an instrument may feel that no further mediation is
needed. The eye was there; the thing was there; what else could be
required? Galileo and Hooke show why the answer is: almost everything.
The eye has to be placed at the instrument. The instrument has to be
made, aligned, and trusted. The specimen or target has to be prepared or
located. The viewing conditions have to be managed. The seen relation
has to be held long enough to be drawn or described. The mark has to be
published in a form that depends on no reader's faith in the original
witness.

Training belongs here too. Instrumented seeing is not automatic the
moment an instrument exists. Observers have to learn how to look, how to
focus, how to ignore irrelevant marks, and how to avoid forcing the view
to match expectation. Training can create reliability, but it can also
create conformity. The record has to carry enough external discipline ---
dates, drawings, comparisons, public demonstrations --- to keep trained
sight open to people outside the training circle.

Replication in this visual setting is more delicate than ``look again.''
Looking again may reproduce the same illusion if the instrument has the
same fault or the observers share the same expectation. The record
therefore needs intermediate objects: sketches, diagrams, plates,
descriptions of the instrument, and instructions for trying again. These
objects give replication a target.

The printed page becomes a witness-management device. It cannot make the
reader see exactly what the observer saw, but it can standardize the
report enough for many readers to argue over the same mark. A drawing of
Jupiter's companions beside dates can be wrong, but if it is wrong it is
wrong in a recoverable way. A plate in \emph{Micrographia} can
exaggerate, stylize, or select, but its public form lets later readers
ask how the image was made. The page turns a fleeting view into a
durable object of suspicion.

Suspicion is not hostile to science here; it is the form public trust
takes when the object has been mediated. The more an observation depends
on a new instrument, the more valuable disciplined suspicion becomes. It
asks whether the same mark appears under changed conditions, whether the
instrument can be checked on familiar objects, whether the observer has
reported failures, and whether rival explanations have been excluded.
Galileo's and Hooke's records cannot remove suspicion; they provide
something specific for suspicion to work on.

Hooke's plates are especially valuable because they hide nothing of their
constructedness. Their scale, detail, and arrangement announce that a
view has been worked into a public artifact. That worked quality is not
a defect; it is what allows the visual claim to circulate. The danger
would be to forget that the artifact is worked. A reader who treats the
plate as a window misses the apparatus in the page. A reader who treats
it as mere art misses the apparatus behind the page. Experimental
realism requires the middle path: a made image answerable to a resistant
object.

Galileo's sketches are less lavish, but their sequence performs a
related job. The small marks near Jupiter need not be beautiful; they
have to be positioned, dated, and compared. Their austerity is part of
their power. The drawing says, in effect, treat this mark not as an image
but as an entry in a relation. The record asks to be read as a
disciplined series rather than as a marvel.

The same distinction governs photographic plates and electronic counters. A plate can dazzle; a counter can click; a
screen can glow. None of these outputs is experimental entry until the
record says how the mark is to be admitted. The eye has been extended;
extension alone produces no public knowledge until the extension has
been given a disciplined inscriptional body.

Visual evidence carries a source history: where the image came from and
what chain of translation made it legible.

Once instrumented and mathematical seeing become public, critics can ask
whether the abstractions used in the record are legitimate. Is the drawn
line an observed line or a geometric convenience? Is the point a body, a
mark, a shorthand, or a calculated position? Is the smooth curve in a
later diagram a fact or a useful fiction? Telescope and microscope
cannot settle those questions; they make them available. Only after
sight has entered a record can the record's abstractions be challenged.

The lesson is not that seeing believes itself; it is that seeing has to
become accountable. Telescope and microscope create new obligations for
record-making, not merely a longer reach. They ask the observer to show
how a private visual event became a public object. The answer may be
imperfect, mediated, and historically contested. That is not a failure
of experimental science; it is the ordinary condition under which
instrumented vision becomes real enough to argue with. Once that
condition is accepted, sight can join the same public record discipline as the
ledger and the table. It becomes a mark with a carrier, a provenance,
and a future critic, not a miracle exempt from record or ordinary
historical doubt.

\begin{phenom}{The Galileo--Hooke Inscription Effect~\cite{galileo1610,hooke1665}}
\label{ph:galileo-hooke-inscription}

\PhStatement
Instrumented sight becomes experimental only when it leaves a portable mark
whose mediation can be inspected.

\PhOrigin
Galileo's telescopic observations and Hooke's microscopic plates each moved
sight through an instrument into a public artifact: dated sketches in one case,
large printed engravings in the other. Their drawings were not neutral copies.
They were disciplined translations from private seeing to shareable record,
and their authority depended on the conditions under which that translation
could be trusted, repeated, criticized, or refused.

\PhObservation
The observer sees through an apparatus, then selects, draws, labels, dates, and
publishes. Each step can preserve or distort the visual event. The resulting
image is neither raw nature nor mere invention; it is a carrier with a history.

\PhConstraint
An instrumented visual claim becomes admissible only when it carries enough
inscriptional detail for others to ask what was seen, how it was seen, what
was selected, what was omitted, and which alternative readings were excluded.
Magnification by itself licenses no fact.

\PhConsequence
Observation becomes public by accepting mediation. The record is not raw
sight; it is sight carried in a durable form. This makes instrumented vision
weaker than private astonishment in one sense, and stronger in the only
sense science can inherit.
\end{phenom}

\section{Abstraction Under Entry}
\episodecoord{1638 / 1734 | Entry :: ADMISSIBLE -> ENCODED}

George Berkeley, writing in 1734 in \emph{The Analyst}, attacked the loose
use of infinitesimals, vanishingly small quantities used in calculus, and
other abstract mathematical entities in the new physics.
He was not mocking science; he was asking when useful symbols had earned
the right to stand in a serious argument. Berkeley belongs here as a
pressure test, not as an enemy. Galileo's physics needed mathematical
abstraction: ideal planes, geometric lines, accelerations, ratios, times,
and simplified bodies. Berkeley's critique asks a question entry cannot
avoid: when has a useful construction become an admissible scientific mark,
and when has the record mistaken a mathematical stand-in for the thing
itself?

Galileo remains a realist case because the abstraction stays tied to
apparatus, measurement, repetition, and correction. An experiment rarely
enters the record as the world in full; it enters through lines, numbers,
diagrams, timings, ratios, and ideal forms that make a public question
answerable. Galileo's inclined planes, falling bodies, and mathematical
arguments are not copies of nature. They prepare nature for entry by
stripping away much of what would otherwise prevent comparison. Friction,
roughness, local imperfection, and bodily mess are not denied because they
are unreal; they are controlled, minimized, or bracketed so that a relation
can be made visible.

Berkeley's pressure is useful because every such bracketing can become a
confusion. A diagram may begin as a disciplined shorthand and end as a
supposed object. A mathematical limit may begin as a method and end as a
hidden entity. A smooth line may begin as a disciplined shorthand for observations and
end as if it had been seen directly. The entry problem is not whether
abstraction is allowed --- without abstraction, experimental science cannot
move. The problem is whether the abstraction declares its status well
enough for later readers to know what kind of mark has entered.

This is the same issue that instrumented sight raised in another form.
Galileo--Hooke made private seeing public by accepting inscription.
Berkeley--Galileo asks what happens when inscription becomes mathematical.
A dot on a page, a line in a diagram, or a ratio in a table is not
identical with the event. It is an encoded object that lets the event be
compared, argued over, and carried. The encoding can be legitimate, but
legitimacy depends on the reader's ability to reconstruct the rule of
translation.

Galileo's achievement was to make motion answerable to mathematical order.
That required an experimental and literary practice, not only equations.
The record had to say what was being measured, how the situation was
simplified, and why the simplification still licensed a claim about moving
bodies. The apparatus delivered no pure mathematical entities. It delivered
timed and arranged marks that could be fitted to a mathematical description.
The fit is the scientific act; it is also the danger point.

Berkeley's critique names that danger without requiring acceptance of
every part of his philosophy. He asks whether the thinker has smuggled in
an unexamined object because the notation works. In the language of entry,
he asks whether \lean{ADMISSIBLE} has been confused with convenient. A
convenience may be indispensable and still need a warning label. If the
record uses a line, point, infinitesimal, average, ideal surface, or
simplified body, the record needs to show how that object functions. Is it
measured? drawn? inferred? calculated? imagined? used as a limiting case?
The answer changes the kind of claim being made.

Abstractions travel better than the events that licensed them. A diagram
can be copied across classrooms for centuries. A formula can detach from
the bench that disciplined it. A curve can move from notebook to textbook
and start looking older, cleaner, and more inevitable than the experiment
itself. The abstraction becomes portable, and that portability is the risk. It
can carry away a result while leaving behind the entry conditions that
made the result responsible.

The realism rule is not to despise abstraction but to refuse anonymous
abstraction. A mathematical object used in an experimental record
needs a role. It may be a coordinate, a model, a limiting case, a
measurement convention, a correction, or an inferential bridge. Each role
licenses different uses. A coordinate may locate a mark without claiming
to be a physical thread in the world. A limiting case may discipline a
calculation without claiming the limit has been physically reached. A
model may expose a relation while omitting known features. What matters is
that the omission be visible enough for criticism.

This is not a retreat from truth into mere bookkeeping; it is how truth
remains answerable to experiment. If Galileo's abstraction is too crude, a
later apparatus can expose the residue. If Berkeley's suspicion is too
sweeping, the success of disciplined abstraction can answer him in
practice. The record needs no fantasy of unmediated fact; it needs a way
to say which mediations were admitted and why. Abstraction becomes
realistic when it remains tied to the work it was built to carry.

The inclined plane is the central apparatus. It is not nature left alone;
it is a prepared world: a surface, a ball, a timing practice, a chosen
angle, and a reason for slowing fall enough to be measured. The ideal
relation is not found lying on the board. It is extracted through
arrangement. That extraction can be excellent science, but only if the
record keeps the arrangement in view. Without the board, the timing, the
angle, and the imperfections, the mathematical relation begins to look
like an unearned given.

Galileo's reported inclined-plane work uses a prepared channel, a rolling
ball, repeated runs, distances marked along the board, and timing by water
collected and weighed. None of those details is the law of fall. They are
the means by which a fast, difficult motion was slowed and divided into
entries. The abstraction becomes credible because the wood, ball, marks,
water, and repetitions make a relation available for inspection.

Failure of an abstraction need not mean the abstraction was foolish; it
may mean the abstraction carried one regime well and another badly. Smooth
surfaces, point masses, perfect clocks, frictionless motion, and ideal
geometries are not lies merely because no workshop can hand them over as
objects. They are admissible when their use is bounded. They become
dangerous when the boundary disappears. A realistic record names not only
what the abstraction permits, but where it is expected to break.

Berkeley's value is to keep that boundary audible. He reminds the reader
that successful manipulation of signs is not the same as an inventory of
reality. Galileo's value is to remind Berkeley's reader that science often
needs disciplined signs in order to let reality answer at all. A mark that
refuses abstraction may remain too local to compare. An abstraction that
refuses its entry conditions may become too smooth to be corrected.

The same tension appears in every later instrument. A photographic plate
turns radiation into a darkened trace. A Geiger counter turns an event
into a click. A radio receiver turns an electromagnetic disturbance into a
coded signal. Each carrier is an abstraction with a material body --- real
evidence, but not the whole of what happened.

The source check for abstraction is a role check: what work is this
mathematical or diagrammatic object performing in the record? Is it
organizing entries, smoothing them, extrapolating beyond them, naming a
limit, or standing in for a material procedure? A later reader cannot
responsibly use the record until that role is recoverable. The demand is
modest and severe at once: not that every abstraction be eliminated, only
that each abstraction appear under a name it can defend.

A diagram can be more honest than a prose claim. A good diagram shows its
artificiality. It tells the reader that a selection has been made: this
axis, this curve, this body, this relation, this omission. The trouble
begins when the diagram's honesty is forgotten and its clean marks are
treated as if they had arrived without construction. Berkeley's suspicion
and Galileo's practice meet there. The construction needs to be clean
enough to reveal a relation and marked enough to remain accountable.

The caution is not anti-mathematical; it is pro-accountability.
Mathematics gives experiment some of its sharpest carriers, but sharpness
cuts both ways. It can reveal a relation crude description would miss, and
it can make an unearned simplification look exact. Entry accepts the tool
while requiring the record to remember the hand that used it, the bench
that constrained it, and the residue it left behind. Without that memory,
elegance outruns entry, and the record floats free from correction.

\begin{phenom}{The Berkeley--Galileo Abstraction Effect~\cite{berkeley1734,galileo1638}}
\label{ph:berkeley-galileo-abstraction}

\PhStatement
An abstraction enters experimental science only when its role in the record is
public enough to be distinguished from the event it helps carry.

\PhOrigin
Galileo's mathematical treatment of motion --- inclined planes, timed runs,
geometric description --- made experimental marks answerable to numerical
and geometric order. Berkeley's critique of infinitesimals and abstract
mathematical entities later pressed on the status of such objects: when
has a useful symbol earned the right to stand in a serious argument? The
episodes are not rivals. Together they expose the entry condition for
abstraction in experimental science.

\PhObservation
Lines, points, ratios, limits, averages, and ideal bodies can make a record
portable and comparable. They can also hide the difference between measured
mark, calculated object, and useful fiction.

\PhConstraint
An experimental record needs to state, or make recoverable, how its
abstractions function. A carrier that works mathematically is not thereby
licensed as an observed object.

\PhConsequence
Realism requires no expulsion of abstraction. It requires disciplined
admission: the abstraction may carry the event only while remaining
answerable to the conditions under which the event entered.
\end{phenom}

\section{Marconi and the Constraint of Silence}
\episodecoord{1901 | Entry :: DISTINGUISHABLE -> ENCODED}

A receiver that records nothing is still inside the experiment when the
apparatus has been made ready. Silence can be a datum if the time window,
operator, receiver, and code are specified. Marconi's 1901 report of the Morse
letter S crossing from Poldhu to Signal Hill is valuable here because the famous
three dots were surrounded by less glamorous conditions: spark transmission,
aerials, tuning, atmospheric noise, operator expectation, and the possibility
of hearing the pattern one had come to hear.

The transatlantic wireless episode was not a clean modern detection experiment.
A signal was reported under difficult conditions, across an astonishing
distance, through an atmosphere that was not a quiet channel. The receiving
arrangement at Newfoundland depended on improvised aerial support, listening
practice, and a known code; the transmitting station in Cornwall depended on
high-voltage spark apparatus and large aerial structures. The historical
glamour can make the event too smooth: transmitter sends, receiver hears,
distance is conquered. The working problem was rougher. A transported mark had
to be separated from noise, and the record had to show why the receiving
arrangement was competent to make that separation.

Signal is not merely arrival. It is arrival under a code. A click, tone, pulse,
or interruption becomes meaningful only because a prior convention lets the
receiver treat one pattern as one mark rather than another. The reported
transatlantic signal was not a little object traveling across the ocean and
landing intact in a box. It was an encoded disturbance made legible by a
transmitting practice, a receiving practice, an expected pattern, and a rule
for deciding whether the pattern had appeared. The code is part of the
apparatus.

This changes the status of absence. In ordinary storytelling, absence is
nothing. In a listening experiment, absence may be an entry. If the apparatus
is ready, the time window is defined, the receiver is competent, and the
detection rule is known, then silence can say that no signal was received under
those conditions. If any of those conditions fails, silence says much less.
The distinction is central. Not hearing is not automatically evidence. Not
hearing while properly listening can become evidence.

The same is true of noise. Static is not simply the enemy of the experiment. It
is the field in which the experiment has to operate. A signal detector without
noise discipline is a rumor amplifier. It may turn expectation into evidence.
It may accept an accidental mark as a message. It may reject a weak arrival as
irrelevant. The record therefore has to preserve more than the positive mark.
It has to preserve the rule by which signal, noise, and operational silence
were separated.

Marconi's episode brings \lean{DISTINGUISHABLE} to the foreground. Earlier
sections asked how a count becomes public, how a table preserves comparison,
and how an image carries instrumented sight. Here the event is less like a
visible object and more like a decision at a threshold. The receiver presents
no moon, pea, or drawn structure. It presents a possible mark in a
field of interference. The first scientific question is whether the mark is
distinguishable enough to enter.

That threshold is not only technical. It is also psychological and social.
Observers may be eager. Sponsors may be waiting. A dramatic success may matter
to reputations, patents, institutions, and newspapers. None of that proves the
signal false. It means that the entry rule has to be stronger than
enthusiasm. The more a result matters, the less it can rely on the
sincerity of the listener. A public record needs to survive the fact that
people wanted the mark to arrive.

The strongest signal record is therefore not the most heroic one. It is the
one that makes false positives and false negatives visible. What could have
produced the same mark without the intended transmission? What weak arrival
would have been missed? What controls or repeated trials would show that the
receiver was not merely responding to local disturbance? What counts as a
failed reception? These questions belong to entry, not to later skepticism.
They decide whether the signal is admissible in the first place.

Encoding makes the situation more powerful and more dangerous. A code lets a
mark travel far beyond the source. It compresses intention into a pattern that
can pass through a channel and be reconstructed elsewhere. But compression
also strips context. The receiver may get the pattern without the local
conditions of the source; the later reader may get the report without the
listening conditions of the receiver. The encoded mark therefore needs a
provenance: source, channel, apparatus, time, threshold, code, and record.

The abstraction problem is plain. A signal is an electrical disturbance made
legible by a code and a receiving practice. It is not the whole event; it is a
carrier that stands for an event under prepared conditions. If the carrier's
status is forgotten, the report becomes magical: a mark arrives because history
says it arrived. If the carrier is treated as mere fiction, the technical
achievement disappears. The signal may enter the record, but only through the
discipline that separated it from noise.

The ordinary craft around the receiver matters as much as the dramatic
distance. Tuning, grounding, antenna arrangement, detector sensitivity,
listening practice, timekeeping, and logging all belong to the entry. A later
reader needs no full romance of the coast, the wind, or the operator's
anticipation. The reader needs enough apparatus detail to know why a
mark was licensed. The romance may motivate memory; the apparatus licenses
evidence.

The negative space of the record is especially important. Times when no
signal was heard, conditions under which reception failed, marks rejected
as static, and uncertainty about what was heard are not dead matter. They
are the frame that lets a positive report mean anything. Entry has
learned to make absence public.

Marconi also shows that transport leaves locality in place. The whole point of
wireless is that a mark can travel, but the traveled mark is always received
somewhere, by a particular apparatus, under particular conditions. Universal
ambition passes through local reception. The public claim depends on the local
receiver being described well enough for others to judge its competence.

The rule is blunt: a signal is a positive mark only inside a disciplined
account of non-marks. Static, silence, missed reception, wrong code, bad
timing, and apparatus failure are not background scenery. They define the
conditions under which a pulse can count as evidence. Without them, a
pulse is too isolated to bear the claim. With them, even a faint mark can become public enough to test, repeat,
and dispute.

Repetition changes the case, but only if repetition is itself disciplined. To
listen again is not enough. The record has to say whether the same code was
expected, whether the same time window was used, whether the apparatus had been
changed, whether the operator knew when to expect the mark, and whether failed
listening was entered beside successful listening. Otherwise repetition may
only repeat expectation. A robust signal history needs a pattern of attempts,
not a single dramatic hearing surrounded by silence that no one bothered to
classify.

This gives a null result its proper dignity. A failed reception under a
competent protocol is not an embarrassment to hide at the edge of the story. It
is information about distance, weather, apparatus, tuning, timing, or the
limits of the detection arrangement. A failed reception outside such a protocol
is weaker; it may say only that no one heard anything useful. The difference
between those two absences is one of the great practical differences in
experimental life. Entry has to mark which kind of absence has occurred.

The code also protects against a certain kind of vagueness. Hearing "something"
is not the same as receiving a signal. The expected pattern gives the receiver
something to test against, but expectation is double-edged. It makes detection
possible and bias possible. The record therefore needs to preserve the relation
between pattern and listener: what was expected, what was heard, how
close the match had to be, and what deviations would have defeated the
claim. A code
without a rejection rule is just a hope with a rhythm.

Modern readers need not choose between easy celebration and easy
debunking. The
episode is historically important because it forced a new kind of entry problem
into view. The useful questions are local and severe. Which marks entered?
Which non-marks were preserved? Which uncertainties remained? Which later
practices had to tighten the threshold? The answer lies in that accounting,
not in the poster version of the event.

The later detector will inherit this discipline. A photographic plate, a
Geiger counter, a cloud chamber, or a trigger system asks the same question in
different material: when is a mark a mark? The wireless receiver makes the
question audible. An encoded event enters only through a competent contrast
class: static, silence, expectation, timing error, local disturbance, instrument
fault. The farther the mark travels, the more carefully the receiving end needs
to say why this mark, here, now, under this code, was allowed to count. Distance
makes the carrier famous, but the entry still happens at the receiver, in the
small discipline of signal, noise, silence, and refusal. That is where the
ocean-crossing claim becomes evidence rather than an exciting story told
afterward by witnesses.

\begin{phenom}{The Marconi Signal Effect~\cite{marconi1901}}
\label{ph:marconi-signal}

\PhStatement
A transported signal enters the record only against a disciplined field of
non-signals.

\PhOrigin
Marconi's transatlantic wireless experiments made signal detection a public
problem of apparatus, code, distance, and threshold. The event was not merely
that a mark arrived, but that a receiving arrangement was treated as competent
to distinguish arrival from absence, signal from static, and code from
expectation.

\PhObservation
The receiver produces marks under noisy conditions. Some marks count as signal;
some are rejected; some silences matter because the apparatus was listening.
The positive report depends on the recorded discipline of these surrounding
non-marks.

\PhConstraint
The record needs to state the detection rule that separates signal, noise,
and operational silence. Without that rule, a transported mark is rumor,
and a silence is merely an unexamined absence.

\PhConsequence
Encoding lets a signal record travel, but only by making absence and error
part of the same discipline as arrival. The signal becomes public when the
record can say not only what was heard, but what would have counted as not
heard.
\end{phenom}

\section{Chaitin: Compression Is Not Possession}
\episodecoord{1970s | Entry :: ENCODED}

Gregory Chaitin, working in the 1970s, helped develop algorithmic
information theory: the study of how short a program can be while still
generating a given string. The theory brings a hard limit into the problem of entry. A
mark can be carried as a string, and a string can be given to a machine,
but the record has not thereby acquired the shortest description for the
string. A finite sequence may be perfectly public as a sequence and still
resist compression. The bits enter the record; the minimal explanation
may not.

The claim has to be stated narrowly. Pick a fixed universal machine
\(U\) --- the chosen reference computer, or fixed programming language,
against which program length is measured. The Kolmogorov complexity of a
finite string \(x\) is defined by
\[
  K_U(x)=\min\{\, |p| : U(p)=x \,\}.
\]
Here \(|p|\) is the length of the program \(p\). This is not a laboratory
reading; it is a formal minimum over descriptions that make the machine
print the string. Change the machine and exact values can shift by a
bounded constant, while the deep phenomenon remains: many strings cannot
be replaced by shorter descriptions without loss. No slogan about the
world being computation follows from this definition. Encoding creates a
public object whose compressibility may be real in the formal sense while
remaining unavailable to direct computation.

That distinction is the Chaitin entry lesson. A string can be printed,
transmitted, stored, checked bit by bit, and copied. It has entered. But
the question ``what is its shortest program?'' belongs to another order of
claim. In general, there is no procedure that will compute exact
Kolmogorov complexity for arbitrary strings. One can prove upper bounds
by giving a program. One can prove lower bounds in special cases. One can
reason about incompressibility in aggregate. But the exact minimal
description is usually not something the record can simply display.

This is not a failure of encoding; it is the maturity of encoding. The
ledger taught that one more can enter. The table taught that many marks
can acquire order. The image taught that sight accepts a carrier. The
signal taught that a transported mark depends on non-marks. Chaitin now
teaches that a carrier can be maximally explicit without being
explanatory. The bits are there. The shortest description, rule, or
program may not be.

The realism issue is delicate. A formal quantity may be well-defined
without being computable in practice or even computable in principle by a
general method. Such a quantity is not thereby unreal; the
boundary the record needs to keep is between existence in a formal system
and possession by an effective procedure. To assume \(K\) is real is to accept the formal proposition
under its conditions; it is not to claim that the experimenter can
compute \(K\) for the string in hand. Entry preserves that difference.

Chaitin is a useful guard against shallow simulation rhetoric. It is
tempting to hear "string," "program," and "machine" and conclude that a
publicly encoded world has been reduced to computation. His work makes
the opposite pressure visible. Encoding can expose irreducibility. A
sequence may be available as data while refusing a shorter effective
account. Algorithmic information theory makes the record more realistic
by refusing the fantasy that every public string arrives with a public
compression.

Tables, visual plates, and radio receivers all show the same caution. A
table can look ordered because it has columns, but the columns cannot
guarantee a short law. A visual plate can carry a trace, but the trace
may not reduce to the anticipated pattern. A radio receiver can record a
sequence of marks, but the sequence may include noise that is not an
encoded message. The existence of a carrier cannot license the existence
of a simple theory. Compression has to be earned.

In ordinary experimental life, this matters whenever a record is too
quickly explained. A curve is fitted, a dataset is modeled, a signal is
decoded, a pattern is named. Some of those acts are legitimate. Some are
premature compression. The Chaitin discipline asks whether the shorter
account actually carries the entries or merely hides their difficulty. A
successful compression lets the record be reconstructed under a rule. A
false compression throws away residue and calls the loss understanding.

A finite record may be complete as a record and incomplete as an
explanation. The absence of a shorter program is not always provable for
the particular case. Failure to find a shorter program is not proof of
randomness; a model that fits some entries is not proof of minimality.
The source discipline therefore keeps the string, the program, the
proof, and the claim of minimality in separate boxes. That separation is
part of encoded realism.

Chaitin's work gives this separation a sharp mathematical form. The
uncomputable is a formally located refusal, not a mystical fog. A
proposition about program size or halting can be meaningful while no
general computation settles it. The record may be forced to distinguish
what can be stated, what can be proved, what can be computed, and what
can merely be bounded. Digital, automated, and model-heavy science needs
that distinction constantly.

Encoding is availability, not ownership. To encode an event is to make it
available for transport, comparison, and computation. The act cannot
guarantee that the event has been explained, compressed, or reduced. A
good record may carry a mark all the way to a machine and still receive
back a refusal. That refusal belongs in the record too.

The dependence on the chosen machine is not an embarrassment to hide; it
is part of the source discipline. A different universal machine may
change the constant overhead of description, just as a different
instrument may change the calibration terms of a measurement. The mature
response is not to pretend the convention is absent; it is to say what
convention is being used and what features survive changes of convention.
Entry requires no view from nowhere, only a declared carrier.

This makes Chaitin surprisingly close to Tycho and Mendel. A table
heading is a convention that lets entries be compared. A universal
machine is a convention that lets descriptions be compared. In both
cases, the convention can be public, powerful, and limited. The table
cannot make its categories natural by printing them. The universal
machine cannot make its complexity values laboratory facts by defining
them. Both give the record a disciplined frame within which public
questions can be asked.

Proof and computation separate here. A proof may show that some strings
are necessarily incompressible, or that no general algorithm can deliver
the exact complexities one might want. That proof is a record of another
kind, but it is not the same as computing the complexity of the
particular string sitting in the dataset. Experimental culture often
blurs this distinction when it moves from theorem to method. Chaitin
forces the blur back into focus: a theorem can license a limit without
handing the laboratory a universal measuring device.

That limit is productive because it prevents encoded records from turning
into a fantasy of total audit. A digital file looks exhaustively explicit
because every bit can be copied. But exhaustiveness of inscription is not
exhaustiveness of understanding. The file may be present without its
rule, law, minimal program, or generative story being present. The
record has become portable, and still the residue remains.

The proposition "K exists" can therefore be treated carefully. Within the
formal setup, \(K\) is not a mood or a metaphor; it is a defined object.
But the object's existence leaves the truth/access distinction intact.
The record may be entitled to speak of \(K\) while being unable to
compute \(K\) exactly for the case it cares about. That is not a paradox
to decorate; it is a discipline: true, perhaps, under the formal
assumption; not therefore available as an entry in the experimental
ledger.

This has consequences for any claim that computation has swallowed
physics. A simulation may produce a record, but the record still has to
enter. The code needs to be specified, the machine needs to be
characterized, the outputs need to be checked, and the compression claim
needs to be distinguished from the encoded string. Chaitin's lesson
neither proves nor disproves a simulated universe; it removes a cheap
step in the argument. Encoding is not the same as explanation, and
existence of a formal quantity is not the same as effective possession by
an observer inside a record.

The realism here is mathematical rather than speculative: algorithmic
information theory keeps accounts at the edge of computation; it names
what has entered, what has merely been bounded, and what remains absent.
What has entered? The string. What else has entered? Perhaps a program,
an upper bound, a proof of a general limit, or a machine convention.
What has not entered? In general, the exact minimal program. That
missing entry is not nothing. It is a structured absence the record
cannot erase.

Clay, tables, drawings, and radio signals may look old-fashioned compared
with computation, but the entry problem has not disappeared. It has sharpened.
Digital encoding lets marks travel with unprecedented fidelity, yet the
question of what the mark licenses remains. A string can be copied
perfectly and still refuse compression. A model can run flawlessly and
still fail to explain the string it outputs. The string survives; the
residue survives with it. Entry has become digital, but it has not
become omniscient, complete, or free from formal refusal. The machine
inherits the record's discipline and its declared limits too.

\begin{phenom}{The Chaitin Compression Effect~\cite{chaitin1974,kolmogorov1965}}
\label{ph:chaitin-compression}

\PhStatement
Encoding can make a mark available to machines without guaranteeing a
computable shortest explanation of the mark.

\PhOrigin
Kolmogorov and Chaitin gave mathematical form to the idea that the
information content of a string can be understood through the length of
its shortest effective description, relative to a chosen universal
machine. Chaitin's work sharpened the relation between program size,
randomness, and incompleteness.

\PhObservation
A finite string may be perfectly public while its exact minimal description
is not effectively obtainable. No general algorithm returns exact Kolmogorov
complexity for arbitrary strings. The record can carry the string without
carrying the shortest program for it.

\PhConstraint
The record needs to distinguish encoded availability from possession by
compression. An upper bound, a model, or a convenient program is not
automatically the minimal account.

\PhConsequence
Entry reaches a formal refusal: some recorded strings can be carried
exactly while resisting exact compression. The refusal is not outside
realism; it belongs to what a realistic encoded record needs to preserve.
\end{phenom}

\section{The Photographic Plate: Trace Without Presence}
\episodecoord{late nineteenth / early twentieth century | Entry :: ENCODED -> ADMISSIBLE}

The photographic plate changes the witness. A person no longer has to be
looking at the moment a mark is made. Light can darken an emulsion, the plate
can be stored, and a later worker can inspect a trace after the event has
passed. This is a magnificent gain for experimental science, but it is not a
magical escape from mediation. The plate is not memory made perfect. It is a
chemical instrument that needs exposure, development, reading, comparison,
and admission.

The plate matters here because it separates occurrence from observation in a
new way. The event may happen while no human eye sees it. The later reader sees
not the event, but a material trace left under prepared conditions. That trace
can be more durable than testimony and more patient than the observer. It can
also be scratched, fogged, overexposed, underexposed, mislabeled, misfiled, or
read under the wrong expectation. Entry has gained a new carrier, and the new
carrier brings its own discipline.

Astronomical plates show the force of the change. A faint star, variable
object, nebular feature, or spectral mark may enter because a plate kept
receiving light after the eye would have tired or failed. The observatory can
become an archive of exposures. Later workers can compare plates across nights,
years, instruments, and catalogues. Henrietta Leavitt used plates to compare
variable stars. Edwin Hubble later relied on plate records of Cepheids
and nebulae to argue for cosmic distance. The plate gave the sky a second, inspectable life.

That second life is not automatic. Exposure time is an entry decision.
Sensitivity of the emulsion is an entry decision. The telescope, guiding,
focus, weather, date, plate handling, development process, and catalogue label
all become part of the mark's provenance. A dot on a plate is not yet a star.
It may be a star, a defect, a dust mark, a scratch, a development artifact, or
a misread trace. The plate preserves a candidate. The record still has to make
the candidate admissible.

The plate gives the same problem a material body. A glass support, a
silver-halide emulsion, a dark slide, a telescope drive, a development bath,
and an archive label all sit between event and later reading. Even a seemingly
mechanical trace has a history. The emulsion is a threshold device, answering
to chemistry, storage, timing, optics, and handling. A trace is an event that
has passed through a prepared surface.

The plate also reshapes the meaning of absence. If no star appears on a
plate, what follows? Perhaps it was too faint. Perhaps the exposure was too
short. Perhaps the plate was damaged. Perhaps clouds intervened. Perhaps the
object was absent or variable. A negative plate becomes evidence only when the
record says enough about the conditions under which the object would have
left a mark. As with Marconi's silence, absence becomes public only under a
competent listening or receiving arrangement.

The strength of the plate is that it invites rereading. A human report may be
frozen in words; a plate can be returned to. Later instruments, better
measuring techniques, and different questions can extract new information from
old exposures. That is why archives matter. The plate can become more useful
than its first interpretation. It may carry a mark whose significance was not
visible to the first worker. Entry has produced a durable object with future
surplus.

That surplus can tempt overconfidence. The plate looks impersonal, and so it
can be mistaken for automatic objectivity. But a plate interprets nothing by
itself. It knows nothing about which marks matter, guarantees nothing about
the catalogue heading, and remembers nothing about why one exposure was
kept and another discarded unless the archive says so. Its impersonality is a
resource, not a license. The plate reduces some human weaknesses and introduces
new material ones.

The public task is therefore to keep plate, trace, and interpretation
separate. The plate is the carrier. The trace is the mark on the carrier. The
interpretation is the claim that the mark stands for an object, event, or
relation. Confusing those levels makes the record brittle. Keeping them apart
makes later correction possible. A disputed dot can be remeasured, compared
with another plate, checked against a catalogue, or rejected as an artifact.

The plate's realism lies in its capacity to hold still while better questions
arrive. It gives the event a material afterlife. The observer leaves. The night
ends. The source changes. The trace remains, if the archive preserves it and
the provenance is strong enough. That afterlife lets science discover that the
original entry was incomplete without losing the mark.

Plate-reading is itself a craft. A trained reader decides where the mark
begins and ends, how strong it is, how it compares with neighboring marks,
and whether the background has been properly understood. Measuring
machines and cataloguing systems can discipline that craft, but they
cannot abolish it. The plate asks
for eyes trained by apparatus and apparatus checked by eyes. It is not a return
to naked seeing. It is a new partnership between stored trace and interpretive
routine.

Comparison is the plate's great public power. A single plate may preserve a
candidate mark; a series of plates can make change visible. Brightness can vary.
Position can shift. A suspected object can appear, fade, or fail to appear
under different exposures. The table returns inside the image archive. Each
plate is a visual entry, but the scientific force often comes from ordered
comparison across entries. The archive is a table made of images.

The plate also introduces a new kind of delayed witness. A worker who was not
present at the telescope may make the decisive comparison later. This is not
secondary in a lesser sense. It is one of the plate's purposes. The original
exposure and the later inspection are different acts. Entry has become
distributed across time: capture, preservation, cataloguing, measurement,
comparison, and interpretation. A record that hides that distribution makes the
plate seem simpler than it is.

The archive can also preserve error. A mislabeled plate may mislead many later
readers. A contaminated emulsion may become a false object. A selection rule
may make an archive overrepresent successes and lose failures. This matters
because archives gain authority by accumulation. A drawer of plates can feel
more objective than a notebook, but accumulated traces are only as trustworthy
as the conditions under which they were admitted and retained. The archive has
an entry rule too.

The plate's relation to time is especially important. It freezes an interval,
not an instant outside time. Exposure duration, guiding stability, and source
motion all affect what the trace means. A streak, blur, or darkness may belong
to the event, to the instrument, or to the exposure practice. The record therefore
needs to tell the reader how the plate integrated time. Without that
information, a trace can be too literal and too vague at once.

By contrast, the Geiger counter turns an event into a discrete mark. The
plate holds a visible trace for later inspection; the counter pins a
moment. One is spatial and archival, the other temporal and punctate.
Both require threshold discipline. Both make events public by refusing
to be simple witnesses. The plate says: here is a lasting surface, but
the record still asks what made the mark. The counter says: here is a
click, but the record still asks what allowed it to count.

A plate is therefore a witness with a body. It has weight, age, chemistry,
scratches, edges, labels, storage conditions, and a route through institutional
hands. Those ordinary facts are not outside the evidence. They are how the
evidence remains tied to the event. A pristine-looking reproduction in a paper
may hide the physical object that carried the original mark. A realistic source
check asks for the body behind the image: which plate, which exposure, which
archive, which handling, which measurement, and which comparison set.

This bodily character is why the plate can discipline imagination. A worker may
hope to find an object, but the plate can refuse. It can show a blank where a
trace was expected, a defect where a discovery was hoped for, or a faint mark
that survives reinspection when enthusiasm fades. The plate is not impartial by
magic. It is resistant because it is material. Its resistance becomes useful
only when the record keeps the material chain available.

The plate thus teaches a quiet lesson about objectivity. Objectivity is not the
absence of hands; hands prepared, exposed, developed, stored, and measured the
plate. Objectivity is the survival of enough material discipline that later
hands can check the earlier ones. The mark becomes public because it is no
longer only what someone remembers seeing. It is something that can be returned
to, handled under rules, compared with neighbors, and, when necessary, refused.
That is why a dark speck can become more public than a brilliant private view:
the speck stayed, and staying made criticism possible after the observing
moment ended and the witness had gone elsewhere.

\begin{phenom}{The Photographic Plate Effect~\cite{leavitt1912,hubble1929}}
\label{ph:photographic-plate}

\PhStatement
A photographic plate makes an event public by giving it a durable trace, but
the trace is admissible only through the plate's material provenance.

\PhOrigin
Astronomical photography made stored exposures central to later measurement,
classification, and comparison. Work such as Leavitt's variable-star studies
and Hubble's distance arguments depended on plates as archived visual records,
not merely on immediate sight.

\PhObservation
The plate can preserve marks no eye witnessed directly. Those marks can be
remeasured, compared, and reinterpreted, but they can also be artifacts of
exposure, development, handling, or catalogue practice.

\PhConstraint
The record needs to distinguish carrier, trace, and interpretation. A dot,
streak, or darkening enters only with enough provenance --- exposure,
development, handling, archive, and catalogue --- to ask what else could
have made the mark.

\PhConsequence
The photographic plate gives observation a material afterlife. That afterlife
makes science more public, but only because the trace remains open to later
inspection, correction, and refusal.
\end{phenom}

\section{The Geiger Counter: A Click Is an Admitted Event}
\episodecoord{early twentieth century | Entry :: DISTINGUISHABLE -> COUNTABLE}

The Geiger counter turns an event into a click, and that simplicity is
exactly what makes it dangerous to describe too quickly. In a typical gas
counter, radiation ionizes gas in a tube; the resulting electrical pulse is
registered as a click or count. A click sounds like a fact. It is discrete,
audible, countable, and public in a way that a faint visual trace may not
be. But the click is not the particle, the decay, or the whole interaction.
It is an admitted event produced by a detector with a threshold, a recovery
time, a background, and a rule for counting.

The counter is a different carrier rather than a better witness. The plate
holds a spatial trace; the counter produces a temporal mark. The plate asks the
reader to inspect a surface. The counter asks the listener, scaler, or register
to accept a pulse as one event. In both cases the apparatus creates evidence by
narrowing the world. The counter narrows especially hard: everything becomes
click or no click.

That narrowing is powerful because it lets events accumulate into rates and
totals. A faint phenomenon that would be almost impossible to hold in memory
can become a series of counted marks. The observer need not draw every
encounter. The instrument increments the record. This is Peano--Kushim
returning inside modern physics: one more, one more, one more, under a rule.
The difference is that the unit is now produced by a detector rather than a
scribe or gardener.

The detector's threshold is the heart of the entry. Too low, and noise enters
as event. Too high, and real events disappear into silence. The record therefore
needs to say what counts as a click, how the apparatus was set, what
background rate was expected, and how the instrument behaves after a count.
Dead time matters because a detector that has just clicked may be unable to
admit another event immediately. A count is not just a sound; it is a sound in
an apparatus state.

This gives the click a disciplined ambiguity. It is more public than a private
flash seen by one observer, but it is also less descriptive. A click says that
the detector admitted an event of the relevant kind. By itself, it says
little about the event's energy, path, source, or cause. The power of the
counter lies in abstraction: the apparatus throws away much of the event in
order to preserve countability. That loss is legitimate only when the record
keeps it visible.

Background is the counter's version of silence. A detector may click even when
the intended source is absent. Cosmic rays, contamination, electronics, and
ambient radioactivity can all contribute. The positive count therefore depends
on a background discipline. The record needs to show how the instrument
behaves without the source, how long the count was taken, and how much
fluctuation is ordinary.
Without that, a total is not a measurement. It is just arithmetic applied to
unclassified clicks.

The counter also trains a different kind of witness. The observer no longer
claims to have seen the event. The observer claims that the apparatus was
competent to admit it. This is a deep shift in experimental culture. Trust
moves from the eye to the detector arrangement, but it never becomes
automatic. Calibration, shielding, geometry, source placement, electronics,
and timing become the new credibility conditions. The apparatus has to earn
the right to click for the record.

The Geiger counter dramatizes the conversion of a hidden occurrence into a
public unit. The unit is wonderfully useful, but it is not innocent. A count
can be wrong because the gas tube discharged on noise, because two events
arrived inside one dead-time interval, because the background was mishandled,
because the source geometry changed, or because the click was counted under
the wrong run. Each failure is an entry failure before it is a theoretical
problem.

Rutherford's alpha-scattering work shows why such discipline matters. The
Manchester experiments relied on alpha particles passing through very thin
metal foil and on scintillations counted on a zinc sulfide screen, often
through a microscope in a darkened room. This was not yet the click of a
Geiger counter; it was the same entry problem in an earlier counting
practice. The theoretical drama of the atom depended on the public handling
of small, rare deflections. A large conclusion could stand only because
individual entries were made countable under a detector and observation
practice.

The counter therefore completes a small arc. The ledger counted goods. The
table counted ordered observations. The plate preserved traces. The radio
receiver distinguished signal from noise. The counter admits hidden events as
units in time. It says: the event may be invisible, but if the apparatus is
disciplined, its entry can still be public.

Rate is where the click becomes more than a tally. Ten clicks in one interval
and ten clicks in another may not carry the same meaning if the intervals,
backgrounds, geometries, or efficiencies differ. A rate is a count with time
attached, and time brings its own entry rule. When the run began, when it
ended, whether the detector stayed stable throughout, whether the source
moved, whether the background was subtracted under comparable conditions
--- each is an entry question. The counter makes numbers
easy to produce, which means it also makes careless numbers easy to produce.

The counter also forces error into the public record. Random arrivals fluctuate.
Small counts are noisy. A single impressive click may mean little; an
accumulated rate may mean much more, but only with uncertainty attached. The
instrument therefore teaches that entry and error are not enemies. Error bars,
background estimates, and repeated runs are ways of making the count
honest. They strengthen the click. They tell later readers how strongly
the click may be used.

This matters for realism because detector evidence is often remote from common
experience. No one sees an atomic nucleus scatter an alpha particle in the
ordinary sense. The detector and its associated record build a public bridge
from invisible event to countable mark. A realist account cannot leap over
that bridge. It needs to inspect the planks: threshold, rate, calibration,
background, geometry, and statistical treatment. The more invisible the
event, the more visible the entry discipline has to become.

The click's austerity is therefore a virtue and a warning. It strips the event
down to admission. That makes large-scale counting possible, but it also means
that later theory cannot pretend the click carried more than it carried. A click
can constrain a theory of matter only when the record says what kind of event
the click was allowed to stand for. The unit is public because it is narrow.
Its narrowness is the price of its strength.

Calibration is the ordinary ritual that keeps the narrowness honest. The
counter needs to be tested against known conditions, checked for drift, and placed
in a geometry that the record can describe. Shielding is not background
decoration; it helps define what the counter is allowed to hear. Source
distance is not a convenience; it changes the relation between event and
admission. Even the instrument's silence during a control interval matters.
Calibration teaches the reader that the click was not merely noticed. It was
licensed.

The final caution is simple. The counter makes entry seem automatic, but
automation is only delegated discipline. Someone still chose the apparatus,
set the threshold, defined the run, handled the background, and interpreted
the total. The machine can reduce the observer's burden, but it cannot
remove the record's obligation to say how a hidden event became one count
in public and under what conditions, limits, and refusals the apparatus
operated.

\begin{phenom}{The Geiger Counter Effect~\cite{rutherford1911}}
\label{ph:geiger-counter}

\PhStatement
A detector click is an experimental entry only when the apparatus has earned
the right to treat the pulse as one event.

\PhOrigin
Early nuclear and radiation experiments depended on making rare, hidden, or
fleeting events countable. Rutherford's alpha-scattering interpretation
depended on disciplined handling of counted marks --- scintillations on a
zinc sulfide screen, not yet the click of a Geiger counter, but already the
same entry problem --- rather than on direct sight of atomic structure.

\PhObservation
The counter converts a hidden interaction into a discrete recorded pulse:
in a gas counter, ionization in the tube becomes an electrical pulse. The
pulse is countable, but it is also abstract --- it preserves admission, not
the full event.

\PhConstraint
The record needs to specify threshold, background, timing, dead time, and
counting conditions. Without those, a click cannot be distinguished from
noise or apparatus behavior.

\PhConsequence
Counting hidden events gives experimental science enormous reach, but only by
making detector discipline part of the fact. The click enters as a licensed
unit, not as the event itself.
\end{phenom}

\section{Faraday's Notebooks: Sequence as Apparatus}
\episodecoord{nineteenth-century laboratory | Entry :: NOTEBOOK -> SEQUENCE}

Michael Faraday was a nineteenth-century experimentalist whose notebooks
recorded the daily craft behind discoveries in electricity and magnetism.
At the Royal Institution in London, where he spent his scientific career
after joining as an assistant in 1813, he worked with little formal
mathematical training by arranging, testing, and rearranging laboratory
apparatus until a relation became visible. His
notebooks --- later published as a multi-volume diary --- run to thousands
of dated entries, sketches, and trial descriptions across decades of
work on electrochemistry, magnetism, and electromagnetic induction.

The apparatus on which Faraday's electromagnetic induction work depended
was deliberately ordinary: insulated wire wound into coils, an iron ring
or core, batteries and switches whose contacts could be made and broken
by hand, and a galvanometer whose needle deflected when current passed.
The decisive observation was a needle deflection that appeared not while
a steady current flowed but at the moments when a circuit was closed or
opened, or when a magnet was moved into or out of a coil. The mark
appeared only when an arrangement changed. A notebook that listed only
the final apparatus and the final result would have hidden where the
experimental work actually happened.

Faraday's notebooks turn a simple claim into a discipline: an experiment
is not public merely because it happened. It becomes available when its
steps, refusals, changes, and repetitions are entered in an order that
another worker can interrogate. The notebook is therefore not a souvenir
of discovery; it is part of the apparatus.

This matters especially for a laboratory where the decisive fact may be a
change in arrangement. Coils are moved. Contacts are made and broken.
Needles deflect, fail to deflect, or deflect in the opposite direction.
A later reader needs more than the final conclusion. The entry needs to
preserve enough sequence to ask what was tried before the successful
configuration appeared and what nearby arrangements failed to register.

The notebook also preserves negative evidence. A failed trial is not just
a personal frustration; it narrows the public field. It says that under
these materials, this geometry, this preparation, and this moment, the
expected mark stayed absent. If such entries are omitted, discovery
becomes theatrical. The record seems to leap from ignorance to result.
Faraday's practice shows a more realistic picture: fact grows through
small entries that include wrong turns, pauses, altered circuits, and
corrections.

Sketches matter for the same reason. A diagram in a notebook is not
decoration. It gives the apparatus a visible grammar. Wires, jars,
magnets, needles, and gaps can be arranged in words, but a sketch catches
relations that prose may blur. The sketch anchors the entry rather than
replacing it. A realistic source check asks whether the published claim
hides notebook labor that fixed the actual geometry of the trial.

The notebook also disciplines memory. Human recollection tends to
compress: first this, then that, therefore the result. A dated page
resists such smoothing. It can show that the important distinction
emerged after repeated failure, that a term changed meaning during the
work, or that an observation was first entered cautiously. That
resistance is not infallibility. A notebook can be selective, hurried,
illegible, or mistaken. But it is a public object in a way memory is not.

Priority also changes under notebook discipline. The question is not
merely who first possessed an idea, but when a claim became tied to a
recoverable procedure. A dated sequence can show whether the decisive
difference was a new material, a reversed connection, a waiting interval,
or a better way to notice motion. It turns invention into a set of
accountable transitions, each small enough to be checked.

Authorship also changes without disappearing. Assistants, instruments,
suppliers, and repeated manual habits can be visible in the page even
when the printed memoir tightens the voice. That visibility is part of
realism. A result is no less real because it required hands, labels, and
bench routines; it becomes more usable when the record shows which hands
and routines mattered.

The notebook also makes replication possible in a way a printed memoir
alone cannot. A later worker who wants to repeat an experiment needs the
practical conditions, not only the abstract conclusion. The notebook can
supply the wire used, the diameter of the coil, the distance to the
magnet, the speed with which contacts were made, the order of trials, and the
apparatus's behavior across the run. Even when the record is silent
about some of these, its silence is informative. A page that omits the
wire gauge invites the later worker to check what gauge was available,
and to ask whether the difference matters.

Faraday's own notebooks document this discipline at the everyday level.
Entries note the date, the materials, the configuration tried, and the
result --- whether the galvanometer needle moved, twitched, or stayed
still. Pages where induction was first noticed sit beside earlier pages
where similar arrangements had given nothing. The famous result depends
on those quieter pages: the deflection that appeared on a particular
morning, when a particular coil was being switched in and out of contact
with a battery, mattered because earlier configurations had been tried
and entered without producing the effect.

The notebook keeps those conditions close to the claim, where later work
can find them and test the published simplicity against the lived
arrangement without pretending the page itself is pure or final.

The notebook makes no claim to truth. It makes an experimental claim
answerable to sequence. It lets later readers separate the reported
effect from the path by which the effect became reportable. Realism
begins where the record keeps the sequence open enough for later
criticism.

\begin{phenom}{The Faraday Notebook Effect~\cite{faraday1932diary}}
\label{ph:faraday-notebooks}

\PhStatement
A laboratory notebook becomes apparatus when it preserves the ordered path by
which a mark, failure, or distinction entered the record.

\PhOrigin
Faraday's experimental notes show discovery as dated work: coils, changing
contacts, and needle deflections appear beside trials, altered
arrangements, sketches, corrections, and observations placed in sequence.

\PhObservation
The notebook gives later readers access to more than a conclusion. It lets
them inspect the succession of attempts by which a conclusion became
admissible: arrangements that produced no deflection, configurations that
produced one, and the differences between them. Each entry preserves the
moment of change --- a contact made, a magnet moved, a coil shifted ---
where the mark could appear.

\PhConstraint
The notebook needs to keep enough detail about arrangement, timing, and
refusal to distinguish experiment from retrospective story.

\PhConsequence
Sequence makes the experiment answerable. The public claim remains tied to
the recorded path that licensed it.
\end{phenom}

\section{Harrison H4: Time Carried as a Record}
\episodecoord{eighteenth-century navigation | Entry :: CLOCK -> POSITION}

John Harrison was an eighteenth-century English clockmaker who spent
decades trying to build a marine timekeeper accurate enough to fix a
ship's longitude. His fourth design --- known as H4 --- was a large
watch-like instrument about thirteen centimetres across, completed in
1759, which carried a reference time accurately enough during voyage
trials in the 1760s to constrain a position at sea. The work was
undertaken in response to the British Parliament's Longitude Prize of
1714, which offered substantial rewards for any method of finding
longitude at sea within stated accuracies.

Longitude is the east-west coordinate of a position on the globe, and it
is intrinsically a time problem. The Earth turns three hundred sixty
degrees in twenty-four hours, which is fifteen degrees of longitude per
hour. If a ship knows that its local noon (when the Sun crosses the
meridian) is happening four hours after noon at a known reference
longitude, the ship is sixty degrees west of that reference. Latitude
can be read from the Sun or stars without a clock. Longitude, by the
chronometer method, needs a carried reference time rather than a direct
reading of position. The chronometer's job is to carry that reference
time, accurately and stably, so the comparison can be made wherever the
ship is.

Harrison's H4 made reference time portable without making time private.
Longitude at sea depends on a comparison: local noon on the ship against
a reference time kept for another meridian. The instrument's task was to
carry one side of the comparison through weather, motion, temperature,
winding, friction, and human handling.

A clock on land can be checked against a familiar routine. A clock at sea
has to hold its rate while the ship moves and the environment changes.
If the rate wanders, the carried time becomes ambiguous. The problem is
disciplined endurance under transport, not simply precision in a cabinet.

The reading looks simple: a hand points, a navigator records. But the
reading is licensed only by the instrument's history. When was it wound?
What was its known rate? How had it behaved before departure? How was it
protected? Was the comparison made under conditions that the record can
describe? A time reading without such questions is like a detector click
without threshold or a plate without provenance.

H4 also shows that workmanship can be part of knowledge. The large
watch-like timekeeper, its jeweled bearings, balance, spring, temperature
compensation, case, and finish were not merely technical ornament. They
helped determine whether the clock could carry time as evidence. The
maker's discipline enters the public result because bad craft would make
the reading untrustworthy. Workmanship can be part of measurement.

The voyage trial is a form of public exposure. The chronometer has to
leave the controlled place where its excellence can be admired and enter
the route where its errors matter. H4 was tested at sea in 1761--62 on a
voyage from Portsmouth to Jamaica, with Harrison's son William carrying
the timekeeper and recording its behaviour during the passage. A second
trial in 1764 took the same instrument to Barbados. In both cases the
chronometer's error at the destination was small enough to fall within
the most stringent accuracy specified by the Longitude Prize. The trials
mattered not only because they succeeded but because they were carried
out under the eyes of skeptical witnesses --- naval officers,
astronomers, and the Board of Longitude --- whose acceptance was part of
the test. The sea is not scenery here. It is the adversarial environment
in which the entry rule is tested.

This example prevents a false opposition between instrument and theory.
The chronometer cannot abolish astronomical reasoning. It gives the
navigator a stable reference against which local observation can be
used. The claim is relational: this local observation, compared with
this carried time, under these rate assumptions, licenses this
longitude. The instrument makes the comparison practical, repeatable,
and contestable.

Even when carried with such care, the chronometer's reading is licensed
by an explicit rate book. Before departure, the instrument's gain or
loss per day was measured against a known standard, often at the Royal
Observatory at Greenwich, a recognised astronomical time standard. The navigator could then apply that rate to convert the
carried time into a corrected reading at each comparison. The rate book
was a living document: checks against celestial events, comparisons with
other timekeepers, and adjustments after handling shocks all entered it.
A longitude derived from H4 was therefore not a bare clock reading; it
was a reading filtered through a documented expectation of drift.

Error is the hinge that makes the chronometer scientific. A wrong rate
not only spoils a voyage; it changes the position the instrument can
license. Thus the useful record is not only the proud final agreement,
but the history of corrections, checks, and disappointments. The record
needs to let drift, adjustment, comparison, and trial count against the
final agreement. Those are not footnotes after the discovery; they are
the means by which the carried mark stops being a private confidence and
becomes an admissible aid.

The navigator then uses time under constraint, as a disciplined companion
to observation. The mark is not fixed to a laboratory bench or an
archive drawer. It travels, and its traveling is part of the claim. A
record that reports only the final longitude hides the fragile middle:
the rate book, the winding routine, the shipboard reading, the
correction, and the comparison. The realist work happens there.

\begin{phenom}{The Harrison Chronometer Effect~\cite{gould1923chronometer}}
\label{ph:harrison-chronometer}

\PhStatement
A chronometer makes longitude enterable when carried time remains disciplined
enough to serve as a public comparison.

\PhOrigin
John Harrison's H4 and its predecessors answered the longitude problem by
carrying reference time through sea conditions, so the ship's local time
could be compared against a known meridian. The 1761--62 Jamaica voyage
and the 1764 Barbados trial showed that the instrument's carried time
could survive a long passage well enough to constrain a position.

\PhObservation
Because the Earth turns fifteen degrees per hour, a time difference
becomes a longitude difference only when the carried time is trusted. The
chronometer reading is therefore not self-sufficient. It depends on rate,
handling, winding, protection, voyage trial, and comparison with local
astronomical observation.

\PhConstraint
The record needs to preserve the conditions under which time was carried
and corrected --- rate, drift, winding routine, voyage trial, and
comparison. Otherwise the reading cannot license position.

\PhConsequence
Portable precision turns navigation into an entry problem: position becomes
public when time survives travel as a checkable mark.
\end{phenom}

\section{Kelvin's Earth: Premise Failure as Entry Refusal}
\episodecoord{nineteenth-century thermodynamics | Entry :: PREMISE -> REFUSAL}

William Thomson, later Lord Kelvin, was one of the nineteenth century's
leading physicists, known for his work in thermodynamics,
electromagnetism, and the engineering of submarine telegraph cables.
Between 1862 and the 1890s, in a series of papers that became central to
Victorian debates about geology and biology, he used thermodynamics to
estimate the age of the Earth.

The argument was disciplined. If the Earth began as a hot body --- whether
molten through and through, or hot at its interior --- and has been
cooling ever since, then the rate at which heat now flows out through its
surface depends on how long the cooling has lasted. Mining and drilling
had established that temperature increases with depth: the upper crust
warms by roughly twenty-five to thirty degrees Celsius per kilometre of
descent. Treat the Earth as a body cooling by conduction, plug in
plausible values for the original temperature and the thermal conductivity
of rock, and the surface-heat-flow observation gives an age. Kelvin's
calculations placed that age between roughly twenty million and a few
hundred million years, depending on the assumptions, with later estimates
tightening toward the lower end. The numbers were transparent, the
physics was the best then available, and the entry conditions were
stated.

The estimate failed because the entry left out two things the Earth
actually has. The first was a source of internal heat that the
nineteenth-century model could not include: the decay of radioactive
elements distributed in the crust and mantle, discovered after 1896.
Radioactive decay supplies heat continuously, so the Earth has cooled
less than a passive body would have over the same span. The second
omission was the way heat moves through the deep interior. Kelvin's
calculation assumed conduction; geophysicists in the twentieth century
established that the mantle convects --- it flows slowly over geological
time, carrying heat from depth to the surface much more efficiently than
conduction alone. Together, internal heat and a different transport regime
changed the cooling problem. The old conduction clock could no longer set
the Earth's age by itself. The modern radiometric age of the Earth is
about four and a half billion years.

Kelvin's calculation shows how a disciplined record can be wrong without
becoming worthless. The calculation treated the Earth as a cooling body
and asked what age would fit the observed thermal gradient under stated
assumptions. Its later failure is no embarrassment to hide; it makes the
entry rule visible.

Kelvin was no fool. The calculation admitted only the sources, boundary
conditions, and interior model available to it. If the Earth had cooled
under those premises, the result would have had force. When later
geology and radioactivity changed the premises, the old number could not
simply be carried forward. The record had licensed a conditional age,
not an unrestricted truth.

This is a different kind of refusal from a blank plate or a silent
receiver. Here the refusal arrives through later knowledge of what the
entry omitted. The calculation was public enough to be inspected; its
assumptions could be named; its constraints could be challenged. That
publicity made correction possible. A vague guess can be forgotten, but a
disciplined wrong calculation can teach the next record where the
missing term entered.

Kelvin also prevents an easy story about theory defeating measurement or
measurement defeating theory. The thermal gradient was a measurement
problem. The cooling model was a physical argument. The age estimate
joined them under premises. When the estimate failed, the failure was
not one side's alone. The apparatus of inference had admitted a world
too simple for the Earth that later evidence required.

A record can be honest and still incomplete. Public discipline is not
omniscience. It is the condition under which incompleteness can be
located. Kelvin's calculation became scientifically useful in part
because later workers could say what it had not included. The mistake
was not pure noise; it was a constrained entry whose constraint became
visible.

This is how failure differs from fraud or rhetoric. Fraud hides the entry
conditions. Rhetoric may evade them. A disciplined failure leaves enough
of itself exposed to be refused. The old age estimate could be rejected
not because it had no structure, but because it had one. The structure
made the wrongness precise.

The later discovery of radioactive heating matters here because it adds
more than a better number. It changes what may be counted as an internal
source. Once that source enters, the old cooling clock is not adjusted
at the margin; its authority is bounded. Later instruments inherit this
lesson. A source term omitted from the entry may remain invisible for
years, yet the eventual correction needs to return to the old
admissibility rule and ask what the earlier record had allowed to exist.

That is not humiliation. It is how a public record learns to carry its
own limits forward, so that later precision can identify what the
earlier record had omitted. Failure, in this sense, is a stored
constraint becoming readable under a larger world. It teaches without
pretending to survive intact forever.

The lesson is sober. Entry gives no guarantee of permanence; it gives a
claim a body in the record: assumptions, marks, calculations, and
admissibility conditions. That body can later be opened. When it is
opened, realism makes no demand that the old claim survive. It demands
only that we know exactly what failed.

\begin{phenom}{The Kelvin Premise-Failure Effect~\cite{thomson1862secular}}
\label{ph:kelvin-earth-age}

\PhStatement
A disciplined calculation can enter the record as a real constraint while
remaining vulnerable to later refusal of its premises.

\PhOrigin
William Thomson, later Lord Kelvin, estimated the Earth's age by treating
surface heat flow and the thermal gradient as evidence of a cooling body
under conduction assumptions, in a series of papers between 1862 and the
1890s.

\PhObservation
The estimate was not a mere guess. It joined measurements and theory
under specific premises --- initial heat, conduction-only transport,
no internal source --- which later evidence made inadequate. Radioactive
heat and mantle convection changed the premises rather than merely
adjusting the arithmetic.

\PhConstraint
The record needs to distinguish the public calculation from the unstated
or unavailable sources it excludes. Otherwise premise failure is mistaken
for simple error.

\PhConsequence
Wrong but disciplined entries are scientifically powerful because they expose
where refusal lands. The record survives as a map of the missing premise.
\end{phenom}

\begin{movementclose}
Entry ends when a mark can survive its moment, but survival is not the same as
truth. A credible realism cannot imagine that the world simply gives facts to a
patient observer. The fact begins as an event, a pressure, a flash, a seed, a
stroke, a click, a time reading, a failed calculation, or an absence. It enters
science only when some apparatus, practice, or record makes it public enough to
be criticized.

That is the ordinary miracle of Entry. Counting is not a child's prelude to
serious theory; it is the first discipline of public order. The Kushim tablet
and Peano's successor function are not historically alike in institution,
purpose, or abstraction, but they rhyme in one crucial way: each refuses a
world where many can remain a blur. The record says one more, then one more,
under a rule. Once that rule exists, disagreement has somewhere to stand.

Tables then teach that entries are not isolated trophies. Tycho's observations
and Mendel's ratios became powerful because individual marks were placed in an
order that survived the observer. A table is a memory with rows, exclusions,
columns, and obligations. It can preserve a pattern, but it can also expose an
odd entry, a bad category, or a premature conclusion. The table's realism is
not that every row is pure. Its realism is that the impurity has a place where
it can be found.

Instrumented sight added another layer. Galileo and Hooke did not merely see
more. They had to teach a public how seeing through an apparatus becomes an
entry rather than a marvel. A lens can disclose, distort, magnify, invert, and
tempt. Berkeley's pressure on abstraction matters here because the record must
not confuse a useful symbol with an unconstrained object. The eye, the drawing,
the instrument, and the mathematical name have to remain distinguishable if
they are to support one another.

Signals make the lesson sharper. Marconi's receiver turns absence into a
discipline: a signal is admitted only against noise, distance, tuning, code,
and the possibility of false pattern. Chaitin's compression warning refuses
another temptation. A short rule may describe a record, but description is not
possession of the world. Complexity can be formal and real inside its formal
setting while still refusing ordinary computation. Entry is strongest when it
admits exactly what its mark can and cannot carry.

The photographic plate and the Geiger counter then split the witness. The
plate gives a trace a body, a surface, a label, a storage history, and a future
reader. The counter gives a hidden event a narrow public unit, a click that can
be counted only under threshold, background, and timing discipline. One stores
a surface for later inspection; the other stores a temporal admission. Both
teach the same severe lesson: a mark is public because it is constrained, not
because it is complete.

Faraday's notebooks, Harrison's chronometer, and Kelvin's failed Earth age
bring that lesson down into ordinary work. The notebook makes sequence
answerable. The chronometer carries time only by carrying rate, correction,
craft, and trial. Kelvin's calculation fails usefully because its premises
remain visible enough to be refused. These are not ornamental historical
episodes. They are the daily grammar of realism: write the path, transport the
mark, expose the premise.

The chapter has therefore not argued that instrumentation manufactures truth.
That would be too easy and too false. It has argued that instrumentation makes
truth claims enterable. Entry is the threshold at which a claim stops being
only experience, hope, habit, memory, authority, or cleverness and becomes a
public object with conditions. Those conditions may be humble: a column heading,
a date, a calibration, a rate, a background run, a catalogue label, a sketch,
or a premise. Humility is not weakness here. It is how the record acquires
teeth.

But every entry also leaves something behind. A count leaves the uncounted.
A table leaves the excluded category. A drawing leaves what the hand simplified.
A signal leaves noise, silence, and undecoded residue. A compressed description
leaves the cost of the machine and the incomputable remainder. A plate leaves
scratches, blanks, and misread specks. A counter leaves dead time and
background. A notebook leaves illegible pauses. A chronometer leaves drift.
A failed calculation leaves the missing source term.

That is where the next movement begins. Once the record can admit marks, the
leftovers become more than inconvenience. They become structured remainders.
Residue is not simply what science has not yet cleaned up. It is what entry
itself makes legible by refusing to absorb everything into the first mark. The
better the entry discipline, the more sharply the leftovers appear. That is why
realism must not rush from entry to conclusion. It must pause over the remainder
and ask what the record has taught us by not fitting.

One last caution matters. Entry is not a ladder we climb once and discard. The
same work returns inside residue, staging, transport, compilation, and harm. A
detector threshold returns inside discovery statistics. A table returns inside
climate synthesis. A notebook returns inside protocol. A chronometer returns
inside synchronization. A failed premise returns inside every model whose
source pass is incomplete. Before asking what a record proves, ask how the
record became a record. That question keeps realism close to practice instead
of letting it float above apparatus, checking, and revision.

So Entry closes with a double result. First, the fact is never merely the event;
it is the event as licensed by a shareable record. Second, the licensed record
is never the whole world; it is the place where the world has become answerable
under specified limits. The next chapter begins with those limits, because
every serious record eventually discovers that its most interesting instruction
may come from what it could not admit on the first pass.
\end{movementclose}
